%
%  chapter_deepdive.tex
%  Aprofundamento Técnico: Por Detrás do Play4Change
%
%  Este capítulo percorre cada interação significativa da plataforma —
%  desde o cliente até à API, serviços de aplicação, agente de IA e base
%  de dados — documentando contratos HTTP, conteúdo de prompts, caminhos
%  de erro e instrumentação de observabilidade.
%
\chapter{Aprofundamento Técnico: Por Detrás do Sistema}
\label{cha:deepdive}

\section{Enquadramento}
\label{sec:deepdive_enquadramento}

Os capítulos anteriores descreveram \emph{o que} o \textit{Play4Change} faz.
O presente capítulo explica \emph{como} o faz — percorrendo cada ação
relevante no portal de administração e na jornada do aprendente desde o
momento em que uma interação é iniciada até ao momento em que os dados
são escritos na base de dados e uma resposta é devolvida ao cliente.

A plataforma serve \textbf{três tipos de cliente} com superfícies
tecnológicas distintas (Figura~\ref{fig:clientes_visao_geral}):

\begin{description}
  \item[\textit{Portal Web} de Administração] implementado em React
    TypeScript com TanStack Query, reservado aos administradores de
    conteúdo; é o único cliente com capacidade de criação de tópicos,
    visualização de progresso SSE em tempo real e gestão de utilizadores.

  \item[Aplicação Móvel] implementada em Kotlin Multiplatform com Compose
    Multiplatform para Android e iOS; é a superfície primária dos
    aprendentes — inscrição em tópicos, resolução de tarefas diárias,
    progressão no mapa de aprendizagem e remediação adaptativa.

  \item[API REST Spring Boot] serve ambos os clientes acima através do
    mesmo conjunto de \textit{endpoints} (com controlo de acesso por
    papel); não existe uma API separada por tipo de cliente.
\end{description}

\begin{figure}[htbp]
  \centering
  \begin{tikzpicture}[>=Stealth, font=\small,
    client/.style={draw, rounded corners=5pt, fill=#1,
                   minimum width=3.2cm, minimum height=1.2cm, align=center},
    srv/.style={draw, rounded corners=5pt, fill=green!12,
                minimum width=4cm, minimum height=1.4cm, align=center},
    db/.style={draw, rounded corners=3pt, fill=gray!12,
               minimum width=2.4cm, minimum height=0.9cm, align=center}
  ]
    \node[client=blue!12]    (web)   at (-4.5, 0) {Portal Web\\{\small React + TypeScript}\\{\tiny Administrador}};
    \node[client=orange!15]  (mob)   at ( 4.5, 0) {App Móvel\\{\small KMP + Compose}\\{\tiny Aprendente}};
    \node[srv]               (api)   at (0,   -2.8){Spring Boot API\\{\small REST + SSE}\\{\tiny HTTPS / JSON}};
    \node[db]                (pg)    at (-2.8,-5.2){PostgreSQL\\{\tiny + pgvector}};
    \node[db]                (minio) at ( 0,  -5.2){MinIO\\{\tiny S3-compatível}};
    \node[db]                (ai)    at ( 2.8,-5.2){Mistral AI\\{\tiny LangChain4j}};

    \draw[->] (web)  -- node[left, font=\tiny]{HTTPS + JWT} (api);
    \draw[->] (mob)  -- node[right,font=\tiny]{HTTPS + JWT} (api);
    \draw[->] (api)  -- (pg);
    \draw[->] (api)  -- (minio);
    \draw[->] (api)  -- (ai);
    \draw[->, dashed, color=blue!60] (api) -- node[above,font=\tiny,sloped]{SSE} (web);
  \end{tikzpicture}
  \caption{Visão geral dos três clientes e dos serviços de suporte do \textit{Play4Change}}
  \label{fig:clientes_visao_geral}
\end{figure}

Para cada fluxo, a análise cobre: o pedido HTTP exato emitido; a sequência
de serviços de aplicação invocados; os dados persistidos; os \textit{prompts}
de IA injetados (quando aplicável); a resposta devolvida no percurso de
sucesso; e o tratamento de erros em cada caminho de falha.
O capítulo cobre explicitamente os dois clientes — portal web e aplicação
móvel — identificando as diferenças de implementação sempre que relevante.

% ─────────────────────────────────────────────────────────────────────────────
\section{Os Três Pilares: Perspetiva de Engenharia}
\label{sec:tres_pilares}

O \textit{Play4Change} está organizado em torno de três pilares de desenho:
\textbf{Envolvimento}, \textbf{Personalização} e \textbf{Reconhecimento}.
Do ponto de vista da engenharia, cada pilar corresponde a um subsistema
distinto.

\begin{description}
  \item[Envolvimento] é concretizado pela camada de gamificação:
    atribuição diária de tarefas (\texttt{TaskAssignment}), pontos de
    experiência, contadores de sequência diária (\textit{streak}) e a
    visualização do mapa de aprendizagem no cliente móvel.
    O servidor impõe uma atribuição por dia por tópico, atribuindo pontos
    apenas na primeira submissão correta.

  \item[Personalização] é a contribuição de investigação central do
    sistema, concretizada por dois mecanismos cooperantes: o
    \textit{pipeline} generativo de IA (LangChain4j + Mistral), que cria
    tarefas específicas por tópico e ramos de remediação adaptativos a
    pedido, e o \textit{motor de similaridade vectorial} (\textit{pgvector}),
    que reutiliza progressivamente ramos gerados anteriormente sempre que
    o perfil de dificuldade de um aprendente corresponde ao de um perfil
    passado.
    Desta forma, o sistema torna-se mais inteligente a cada aprendente que
    serve: um padrão de dificuldade registado por um utilizador enriquece
    a biblioteca de ramos reutilizáveis disponíveis para todos os
    utilizadores futuros.

  \item[Reconhecimento] é implementado pelo subsistema de
    \textit{badges} (\texttt{BadgeIssuanceService}), que avalia critérios
    configuráveis — conclusão de módulo, comprimento de sequência diária,
    limiares de pontuação — após cada submissão de tarefa, emitindo
    \textit{badges} digitais verificáveis armazenados em PostgreSQL.
\end{description}

O restante do capítulo organiza-se em torno destes pilares, começando
pela fundação de autenticação comum a ambos os tipos de ator.

% ─────────────────────────────────────────────────────────────────────────────
\section{Autenticação: O Protocolo \textit{Magic Link}}
\label{sec:auth_deepdive}

O \textit{Play4Change} adopta uma estratégia de autenticação sem palavra-passe
(\textit{magic link}), eliminando o armazenamento de credenciais e a
superfície de ataque associada a ataques de reutilização de
credenciais~\cite{owasp2021a07}.

\subsection*{Passo 1 — Pedido de ligação}

O administrador ou aprendente submete o seu endereço de correio eletrónico.
O \textit{frontend} web envia:

\begin{verbatim}
POST /auth/magic-link
Content-Type: application/json

{ "email": "alice@example.com" }
\end{verbatim}

O \texttt{AuthController} delega no \texttt{MagicLinkService}, que:
(i)~localiza ou cria o registo de utilizador;
(ii)~gera um \textit{token} aleatório criptograficamente seguro e armazena
o seu resumo SHA-256 com um TTL de 15 minutos;
(iii)~envia o \textit{e-mail} com a ligação de autenticação via API Resend.
A resposta é \texttt{202 Accepted} com uma mensagem genérica — a mesma
resposta é devolvida independentemente de o endereço existir ou não,
impedindo ataques de enumeração de utilizadores~\cite{owasp2021a05}.
Um \texttt{RateLimitService} impõe um máximo de três pedidos por endereço
por hora.

\subsection*{Passo 2 — Clique na ligação}

O \textit{e-mail} contém uma ligação para \texttt{GET /auth/verify?token=\{raw\}}.
O servidor \emph{não} consome o \textit{token} neste passo; em vez disso,
emite:

\begin{verbatim}
HTTP/1.1 302 Found
Location: play4change://auth/verify?token={raw}
\end{verbatim}

No cliente móvel, o sistema operativo interceta o esquema URI personalizado
e abre a aplicação nativa.
No portal \textit{web}, o \textit{frontend} interceta o redirecionamento
através da página de \textit{callback} OAuth (\texttt{AuthCallbackPage}).

\subsection*{Passo 3 — Verificação do \textit{token}}

O cliente envia o \textit{token} em claro:

\begin{verbatim}
POST /auth/magic-link/verify
Content-Type: application/json

{ "token": "{raw}" }
\end{verbatim}

O \texttt{MagicLinkService} recalcula o resumo SHA-256 do \textit{token}
recebido, consulta o \texttt{MagicLinkTokenJpaRepository} em busca de um
registo correspondente não expirado e, em caso de sucesso, marca o
\textit{token} como consumido (utilização única).
O \texttt{TokenService} emite então um \textit{token} de acesso JWT
(curta duração, 15--30 minutos) e um \textit{token} de atualização
(TTL de 7 dias) armazenado em PostgreSQL.
A resposta é:

\begin{verbatim}
HTTP/1.1 200 OK
{
  "accessToken": "<JWT>",
  "refreshToken": "<opaque>",
  "expiresIn": 900
}
\end{verbatim}

No \textit{frontend} web, o \texttt{apiClient.ts} armazena o \textit{token}
de acesso em \texttt{sessionStorage} e o \textit{token} de atualização
num \textit{cookie} com os atributos
\texttt{SameSite=Strict; Secure}, mitigando ataques
CSRF~\cite{owasp2021a01,west2016samesite}.

Na aplicação móvel, a interface \texttt{TokenStorage} abstrai o mecanismo
de persistência segura com implementações distintas por plataforma:

\begin{itemize}
  \item \textbf{Android} — \texttt{EncryptedSharedPreferencesTokenStorage}
    utiliza a API \texttt{EncryptedSharedPreferences} do Jetpack Security,
    que aplica encriptação AES-256-GCM às chaves e AES-256-SIV aos valores,
    com a chave-mestra armazenada no \textit{Android Keystore}.

  \item \textbf{iOS} — \texttt{KeychainTokenStorage} acede diretamente à
    \textit{Security Framework} do sistema via \textit{Kotlin/Native CInterop}.
    A classe usa \texttt{kSecUseDataProtectionKeychain = true} (obrigatório
    em iOS 13+ e no simulador iOS 26.2) e
    \texttt{kSecAttrAccessibleAfterFirstUnlock}, que permite a renovação de
    \textit{tokens} em segundo plano após o primeiro desbloqueio do
    dispositivo pós-arranque.
    A estratégia de escrita aplica \texttt{SecItemAdd} e, em caso de
    \texttt{errSecDuplicateItem} (reinstalação da aplicação), recorre a
    \texttt{SecItemUpdate}, evitando que uma reinstalação provoque erros
    silenciosos de persistência.
\end{itemize}

\subsection*{Passo 4 — Renovação silenciosa}

\paragraph{Portal web.}
Os interceptores de resposta Axios detetam \texttt{401 Unauthorized}
em pedidos que não sejam de autenticação.
O interceptor emite \texttt{POST /auth/refresh} usando o \textit{token}
de atualização armazenado.
Os pedidos concorrentes que chegam durante a renovação são colocados em
fila através de um padrão de subscritor; assim que o novo par de
\textit{tokens} é obtido, todos os pedidos em fila são reenviados com o
novo \textit{token} de acesso.
Se a renovação falhar, \texttt{clearTokens()} apaga todas as credenciais
armazenadas e dispara um \texttt{CustomEvent("auth:session-expired")} que
a árvore React trata via \textit{React Router}, preservando o encaminhamento
SPA sem recarregar a página.

\paragraph{Aplicação móvel.}
O \texttt{HttpClientFactory} instala o \textit{plugin} \texttt{Auth} do
Ktor com a estratégia \texttt{bearer}.
O bloco \texttt{loadTokens} carrega o par de \textit{tokens} a partir do
\texttt{TokenStorage} no início de cada pedido.
O bloco \texttt{refreshTokens} é invocado automaticamente pelo Ktor quando
o servidor devolve \texttt{401}: envia \texttt{POST /auth/refresh} com o
\textit{token} de atualização e armazena o novo par.
Erros de rede transitórios durante a renovação não limpam os \textit{tokens}
— o pedido original falha e o utilizador pode tentar novamente.
Se a renovação falhar com um erro de servidor (não transitório), os
\textit{tokens} são apagados e o \texttt{SessionEventBus.sessionExpired()}
é emitido num \texttt{MutableSharedFlow}.
O \texttt{DefaultRootComponent} observa este fluxo e invoca
\texttt{navigation.replaceAll(Config.Login)}, redirecionando o utilizador
para o ecrã de autenticação sem destruir e recriar a \textit{Activity}
Android.

\begin{figure}[htbp]
  \centering
  \begin{tikzpicture}[>=Stealth, font=\small,
    client/.style={draw, fill=blue!10, rounded corners=3pt,
                   minimum width=2.8cm, minimum height=0.7cm, align=center},
    server/.style={draw, fill=green!10, rounded corners=3pt,
                   minimum width=2.8cm, minimum height=0.7cm, align=center},
    ext/.style={draw, fill=orange!10, rounded corners=3pt,
                minimum width=2.8cm, minimum height=0.7cm, align=center}
  ]
    \node[client] (browser) at (0,0)    {Browser / Aplicação};
    \node[server] (api)     at (5.5,0)  {Spring Boot API};
    \node[ext]    (email)   at (10.5,0) {Resend (E-mail)};
    \node[ext]    (db)      at (10.5,-5){PostgreSQL};

    \draw[dashed] (0,-0.35) -- (0,-9.5);
    \draw[dashed] (5.5,-0.35) -- (5.5,-9.5);
    \draw[dashed] (10.5,-0.35) -- (10.5,-9.5);

    \draw[->] (0,-1)   -- node[above,font=\tiny]{POST /auth/magic-link} (5.5,-1);
    \draw[->] (5.5,-1.7) -- node[above,font=\tiny]{enviar e-mail} (10.5,-1.7);
    \draw[->] (5.5,-2.4) -- node[above,font=\tiny]{guardar token SHA-256} (10.5,-2.4);
    \draw[->] (5.5,-3.1) -- node[above,font=\tiny]{202 Accepted} (0,-3.1);
    \draw[->] (0,-4)   -- node[above,font=\tiny]{GET /auth/verify?token=X} (5.5,-4);
    \draw[->] (5.5,-4.7) -- node[above,font=\tiny]{302 → play4change://} (0,-4.7);
    \draw[->] (0,-5.7) -- node[above,font=\tiny]{POST /auth/magic-link/verify} (5.5,-5.7);
    \draw[->] (5.5,-6.4) -- node[above,font=\tiny]{verificar SHA-256} (10.5,-6.4);
    \draw[->] (10.5,-7.1) -- node[above,font=\tiny]{marcar consumido} (5.5,-7.1);
    \draw[->] (5.5,-7.8) -- node[above,font=\tiny]{200 OK \{accessToken, refreshToken\}} (0,-7.8);
    \draw[->] (0,-9)   -- node[above,font=\tiny]{Pedido API (+ Bearer token)} (5.5,-9);
    \draw[->] (5.5,-9.5) -- node[above,font=\tiny]{200 OK / 401 → renovação} (0,-9.5);
  \end{tikzpicture}
  \caption{Sequência de autenticação por \textit{magic link}}
  \label{fig:auth_sequence}
\end{figure}

\subsection*{Tratamento do \textit{deep link} na aplicação móvel}

Quando o utilizador toca na ligação do \textit{e-mail} num dispositivo
móvel, o sistema operativo interceta o esquema URI
\texttt{play4change://auth/verify?token=\{raw\}} e abre a aplicação.

No Android, a \texttt{MainActivity} recebe o \texttt{Intent} no método
\texttt{onCreate()} (arranque frio) ou \texttt{onNewIntent()} (aplicação
já em execução).
O método \texttt{handleIntent()} extrai o parâmetro \texttt{token} da
\textit{query string} quando o caminho começa em \texttt{/auth/verify}:

\begin{verbatim}
val token = intent.data
    ?.takeIf { it.path?.startsWith("/auth/verify") == true }
    ?.getQueryParameter("token") ?: return
root.handleDeepLink(token)
\end{verbatim}

O \texttt{DefaultRootComponent.handleDeepLink()} verifica o \textit{token}
com \texttt{POST /auth/magic-link/verify} e navega para
\texttt{Config.Home} em caso de sucesso.
Se a verificação falhar \emph{e} não existir já uma sessão ativa
(i.e., \texttt{tokenStorage.getAccessToken() == null}), navega para
\texttt{Config.Login} — caso contrário não interrompe a sessão existente,
cobrindo o cenário em que o utilizador toca na ligação num segundo
dispositivo enquanto já está autenticado noutro.

% ─────────────────────────────────────────────────────────────────────────────
\section{Arquitetura do Cliente Móvel: Kotlin Multiplatform}
\label{sec:kmp_arquitetura}

A aplicação móvel do \textit{Play4Change} é construída com
\textbf{Kotlin Multiplatform} (KMP), partilhando lógica de negócio,
camada de rede, módulos de injeção de dependências e código de interface
de utilizador entre Android e iOS através de um único conjunto de fontes
\texttt{commonMain}.
Apenas os componentes que dependem de APIs nativas não abstraídas
permanecem em conjuntos de fontes específicos por plataforma
(Figura~\ref{fig:kmp_estrutura}).

\begin{figure}[htbp]
  \centering
  \begin{tikzpicture}[>=Stealth, font=\small,
    common/.style={draw, fill=blue!10, rounded corners=4pt,
                   minimum width=10cm, minimum height=0.85cm,
                   align=center, text width=9.5cm},
    platform/.style={draw, fill=#1, rounded corners=4pt,
                     minimum width=4.5cm, minimum height=0.85cm, align=center},
    lbl/.style={font=\bfseries\small}
  ]
    % commonMain layers
    \node[common] (ui)   at (0,  0)   {\textbf{commonMain — Interface (Compose Multiplatform)}\\
                                        Ecrãs, componentes de design, animações};
    \node[common] (pres) at (0,-1.1)  {\textbf{commonMain — Apresentação (Decompose + MVI)}\\
                                        \texttt{DefaultRootComponent}, \texttt{BaseComponent}, efeitos, eventos};
    \node[common] (dom)  at (0,-2.2)  {\textbf{commonMain — Domínio}\\
                                        repositórios (interface), modelos, \texttt{AppError}, \texttt{NetworkError}};
    \node[common] (data) at (0,-3.3)  {\textbf{commonMain — Dados}\\
                                        \texttt{HttpClientFactory}, \texttt{TokenStorage} (interface), DTOs, Ktor};
    \node[common] (di)   at (0,-4.4)  {\textbf{commonMain — Injeção (Koin)}\\
                                        \texttt{AppModule}, módulos de funcionalidade, \texttt{platformModule} (expect)};

    % Platform-specific
    \node[platform=green!15]  (and) at (-2.6,-5.8) {\textbf{androidMain}\\
                                        OkHttp, \texttt{EncryptedSharedPreferences},\\
                                        \texttt{MainActivity}, NetworkDetection};
    \node[platform=gray!18]   (ios) at ( 2.6,-5.8) {\textbf{iosMain}\\
                                        Darwin, \texttt{KeychainTokenStorage},\\
                                        \texttt{MainViewController}, NetworkDetection};

    \draw[->] (ui)   -- (pres);
    \draw[->] (pres) -- (dom);
    \draw[->] (dom)  -- (data);
    \draw[->] (data) -- (di);
    \draw[->] (di.south) ++(-1.5,0) -- (and.north);
    \draw[->] (di.south) ++( 1.5,0) -- (ios.north);
  \end{tikzpicture}
  \caption{Estrutura de conjuntos de fontes KMP da aplicação \textit{Play4Change}}
  \label{fig:kmp_estrutura}
\end{figure}

\subsection*{Padrão \textit{expect}/\textit{actual}}

O KMP compila o código \texttt{commonMain} para cada plataforma alvo.
As declarações \texttt{expect} no código partilhado definem contratos que
as implementações \texttt{actual} em cada plataforma devem satisfazer.
Os três pontos de variabilidade no \textit{Play4Change} são:

\begin{description}
  \item[\texttt{TokenStorage}] Interface partilhada; implementada por
    \texttt{EncryptedSharedPreferencesTokenStorage} (Android) e
    \texttt{KeychainTokenStorage} (iOS).
    O módulo Koin da plataforma (\texttt{platformModule}) regista a
    implementação correta em tempo de compilação — o código partilhado
    nunca depende de uma implementação concreta.

  \item[\texttt{HttpEngineFactory}] Declaração \texttt{expect fun
    platformHttpClient(block)}: devolve um \texttt{HttpClient} com motor
    OkHttp no Android e Darwin (NSURLSession) no iOS.
    O motor Darwin é nativo em iOS e não requer dependência de terceiros.

  \item[\texttt{NetworkDetection}] Verifica a disponibilidade de rede
    antes de cada pedido: \texttt{ConnectivityManager} no Android,
    \texttt{NWPathMonitor} no iOS.
    Um \texttt{NetworkError.NoConnection} é emitido sem realizar o pedido
    HTTP, poupando tempo de \textit{timeout}.
\end{description}

\subsection*{Injeção de dependências com Koin}

A injeção de dependências é gerida pelo Koin.
O grafo de módulos divide-se em dois níveis: o \texttt{AppModule}
partilhado (que cria o \texttt{HttpClient}, regista os repositórios HTTP
e configura o \texttt{SessionEventBus}) e o \texttt{platformModule}
(declaração \texttt{expect} no \texttt{commonMain}, implementada por cada
plataforma para fornecer a \texttt{TokenStorage} e o \texttt{NetworkConfig}
com o URL base correto).
Módulos de funcionalidade (\texttt{AuthModule}, \texttt{TaskModule},
\texttt{StruggleModule}, etc.) registam os seus componentes de apresentação
em âmbito \textit{factory}, garantindo que cada instância de ecrã recebe
o seu próprio \texttt{DefaultXxxComponent}.

\begin{table}[ht]
  \caption{Módulos Koin e responsabilidades na aplicação móvel}
  \label{tab:koin_modules}
  \vskip0.5em
  \centering
  \begin{tabular}{lp{8cm}}
    \toprule
    \textbf{Módulo} & \textbf{Responsabilidade} \\
    \midrule
    \texttt{platformModule} (actual) &
      \texttt{NetworkConfig} (URL base), \texttt{TokenStorage} (plataforma) \\
    \texttt{CoreModule} &
      \texttt{HttpClient} (via \texttt{HttpClientFactory}),
      \texttt{SessionEventBus} \\
    \texttt{AuthModule} &
      \texttt{AuthRepository}, \texttt{DefaultLoginComponent} \\
    \texttt{HomeModule} &
      \texttt{HomeRepository}, \texttt{DefaultHomeComponent} \\
    \texttt{TaskModule} &
      \texttt{TaskRepository}, \texttt{DefaultTaskComponent} \\
    \texttt{StruggleModule} &
      \texttt{StruggleRepository}, \texttt{DefaultStruggleComponent} \\
    \texttt{ExploreModule} &
      \texttt{ExploreRepository} (lista de tópicos + inscrição) \\
    \texttt{ProfileModule} &
      \texttt{ProfileRepository}, \texttt{DefaultProfileComponent} \\
    \bottomrule
  \end{tabular}
\end{table}

% ─────────────────────────────────────────────────────────────────────────────
\section{Navegação no Cliente Móvel: Decompose}
\label{sec:decompose_navegacao}

A navegação na aplicação móvel é gerida pela biblioteca
\textbf{Decompose}, que modela a pilha de ecrãs como uma máquina de
estados serializada — compatível com o ciclo de vida Android e com a
destruição de ecrã em iOS.

\subsection*{\texttt{DefaultRootComponent}}

O \texttt{DefaultRootComponent} implementa a interface \texttt{RootComponent}
e serve como único ponto de orquestração de navegação.
Mantém uma \texttt{StackNavigation<Config>} que representa a pilha de ecrãs.
A \texttt{Config} é uma interface selada serializada com oito alternativas:

\begin{verbatim}
sealed interface Config {
    Splash | Login | Home
    Task(userTaskId: String)
    Profile | About | Explore
    Struggle(enrollmentId: String)
}
\end{verbatim}

A serialização \texttt{kotlinx.serialization} permite que o Decompose
persista e restaure o estado da pilha de navegação através de alterações de
configuração Android (rotação do ecrã, modo multi-janela) sem necessidade
de \texttt{ViewModel} ou \texttt{SavedStateHandle}.

\subsection*{Animações de transição}

O \texttt{App.kt} envolve a pilha em \texttt{AnimatedContent} com uma
especificação de transição personalizada baseada em \texttt{CubicBezierEasing}.
A direção da animação é determinada pela comparação do tamanho da pilha:
um aumento de tamanho (\textit{push}) desliza da direita para a esquerda;
uma diminuição (\textit{pop}) inverte a direção.
Esta transição é idêntica nas duas plataformas — o Compose Multiplatform
renderiza-a nativamente no motor de compositor de cada sistema operativo.

\subsection*{Fluxo de ecrãs}

\begin{figure}[htbp]
  \centering
  \begin{tikzpicture}[>=Stealth, font=\small,
    scr/.style={draw, fill=#1, rounded corners=5pt,
                minimum width=2.8cm, minimum height=0.9cm, align=center},
    ev/.style={font=\tiny, align=center}
  ]
    % Nodes
    \node[scr=yellow!20]  (splash)   at (0,    0)   {Splash};
    \node[scr=blue!10]    (login)    at (-3.2,-1.8) {Login};
    \node[scr=green!15]   (home)     at ( 0,  -1.8) {Home};
    \node[scr=orange!15]  (task)     at (-3.2,-3.8) {Task};
    \node[scr=red!10]     (struggle) at ( 0,  -3.8) {Struggle};
    \node[scr=gray!15]    (explore)  at ( 3.2,-1.8) {Explore};
    \node[scr=gray!15]    (profile)  at ( 3.2,-3.8) {Profile};

    % Transitions
    \draw[->] (splash) -- node[left,ev]{sem sessão} (login);
    \draw[->] (splash) -- node[right,ev]{sessão ativa} (home);
    \draw[->, dashed, color=red] (login)  -- node[above,ev]{link verificado} (home);
    \draw[->] (home)   -- node[left,ev]{tarefa} (task);
    \draw[->] (home)   -- node[right,ev]{explorar} (explore);
    \draw[->] (home)   -- node[right,ev]{perfil} (profile);
    \draw[->] (task)   -- node[left,ev]{dificuldade\\detetada} (struggle);
    \draw[->] (struggle) -- node[above,ev]{concluído} (home);
    \draw[->, dashed, color=red] (home) to[bend right=40]
              node[right,ev]{sessão expirada} (login);
  \end{tikzpicture}
  \caption{Grafo de ecrãs da aplicação móvel gerido pelo Decompose}
  \label{fig:decompose_nav}
\end{figure}

A transição \texttt{Splash → Login} ou \texttt{Splash → Home} é decidida
pelo \texttt{DefaultSplashComponent}: se o \texttt{TokenStorage} contiver
um \textit{token} de acesso válido, o efeito \texttt{SplashEffect.NavigateToHome}
é emitido; caso contrário, \texttt{SplashEffect.NavigateToLogin}.
A transição de expiração de sessão (linha a tracejado vermelho) é
despoletada pelo \texttt{SessionEventBus} independentemente do ecrã ativo
— o \texttt{DefaultRootComponent} observa o fluxo num
\texttt{CoroutineScope(Dispatchers.Main + SupervisorJob())} e invoca
\texttt{navigation.replaceAll(Config.Login)}.

\subsection*{Padrão MVI no cliente móvel}

Cada funcionalidade segue o padrão \textbf{MVI}
(\textit{Model–View–Intent}) com três tipos de mensagem:

\begin{description}
  \item[\textit{Events}] Intenções do utilizador (clique num botão,
    submissão de resposta) enviadas do ecrã para o componente via
    \texttt{onEvent(E)}.
  \item[\textit{State}] Estado imutável do ecrã exposto como
    \texttt{Value<S>} (do Decompose), subscrito pelo ecrã via
    \texttt{subscribeAsState()}. Inclui sempre \texttt{isLoading: Boolean}
    e \texttt{error: AppError?}, geridos pelo \texttt{BaseView} de forma
    uniforme.
  \item[\textit{Effects}] Efeitos de navegação de uso único (e.g.,
    \texttt{NavigateToStruggle(enrollmentId)}) emitidos via
    \texttt{SharedFlow} e tratados pelo \texttt{App.kt} com
    \texttt{LaunchedEffect(component)}.
\end{description}

\begin{figure}[htbp]
  \centering
  \begin{tikzpicture}[>=Stealth, font=\small,
    box/.style={draw, rounded corners=4pt, fill=#1,
                minimum width=3.0cm, minimum height=0.9cm, align=center}
  ]
    \node[box=blue!10]    (view)  at (0,    0) {Ecrã Compose\\(\textit{View})};
    \node[box=green!12]   (comp)  at (5.5,  0) {Componente\\(\textit{Presenter})};
    \node[box=orange!12]  (repo)  at (5.5, -2.5){Repositório HTTP\\(\textit{Model})};
    \node[box=gray!12]    (state) at (0,   -2.5){\texttt{Value<State>}\\(imutável)};

    \draw[->, thick]              (view)  -- node[above,font=\tiny]{\textit{Events} (onEvent)} (comp);
    \draw[->, thick]              (comp)  -- node[right,font=\tiny]{chamadas suspend} (repo);
    \draw[->, thick]              (repo)  -- node[below,font=\tiny]{Result / exceção} (comp);
    \draw[->, thick]              (comp)  -- node[left, font=\tiny]{updateState\{\}} (state);
    \draw[->, thick]              (state) -- node[left, font=\tiny]{subscribeAsState()} (view);
    \draw[->, dashed, color=red!60] (comp) to[bend left=25]
            node[above,font=\tiny]{\textit{Effects} (navigate)} (view);
  \end{tikzpicture}
  \caption{Fluxo de dados MVI na camada de apresentação do cliente móvel}
  \label{fig:mvi_fluxo}
\end{figure}

% ─────────────────────────────────────────────────────────────────────────────
\section{O \textit{Pipeline} de Pedidos do Servidor Spring Boot}
\label{sec:pipeline_spring}

Antes de analisar fluxos funcionais específicos, é fundamental compreender
o caminho que \emph{qualquer} pedido HTTP percorre desde a chegada ao
servidor até à emissão da resposta.
O \textit{Play4Change} utiliza o Spring Boot com o servidor \textit{embedded}
Apache Tomcat.
Um pedido atravessa oito camadas distintas em sequência antes de atingir
o controlador de domínio (Figura~\ref{fig:pipeline_spring}).

\begin{figure}[htbp]
  \centering
  \begin{tikzpicture}[>=Stealth, font=\small,
    layer/.style={draw, fill=#1, rounded corners=4pt,
                  minimum width=10.5cm, minimum height=0.8cm, align=center,
                  text width=10cm},
    arrow/.style={->, thick, color=gray!60}
  ]
    % Draw layers top-to-bottom (request direction)
    \node[layer=orange!15]  (nginx)   at (0, 0.0) {%
      \textbf{1. Nginx} — proxy inverso, terminação TLS, desativação de \textit{buffering} SSE};
    \node[layer=yellow!20]  (tomcat)  at (0,-1.1) {%
      \textbf{2. Tomcat} (\textit{embedded}) — aceitação TCP, gestão de fios de execução};
    \node[layer=red!12]     (ratelim) at (0,-2.2) {%
      \textbf{3. \texttt{RateLimitFilter}} \small[\texttt{@Order(HIGHEST\_PRECEDENCE)}]
      — limita pedidos \texttt{/auth/**} por IP};
    \node[layer=red!8]      (reqid)   at (0,-3.3) {%
      \textbf{4. \texttt{RequestIdFilter}} \small[\texttt{@Order(1)}]
      — gera \texttt{X-Request-Id}, injeta no MDC};
    \node[layer=blue!10]    (cors)    at (0,-4.4) {%
      \textbf{5. Spring Security — \texttt{CorsFilter}}
      — valida origem, emite cabeçalhos \texttt{Access-Control-*}};
    \node[layer=blue!15]    (jwt)     at (0,-5.5) {%
      \textbf{6. \texttt{JwtAuthFilter}} \small[\texttt{OncePerRequestFilter}]
      — analisa \texttt{Bearer}, povoa \texttt{SecurityContextHolder}};
    \node[layer=blue!20]    (authz)   at (0,-6.6) {%
      \textbf{7. Spring Security — Autorização}
      — matéria de rotas: \texttt{permitAll} / \texttt{hasRole("ADMIN")} / \texttt{authenticated()}};
    \node[layer=green!12]   (metrics) at (0,-7.7) {%
      \textbf{8. \texttt{WebMvcMetricsFilter}} (auto-configurado)
      — regista \texttt{http\_server\_requests\_seconds\_*} no Prometheus};
    \node[layer=purple!10]  (disp)    at (0,-8.8) {%
      \textbf{9. \texttt{DispatcherServlet}} — mapeamento de controlador, negociação de conteúdo};
    \node[layer=teal!10]    (logpre)  at (0,-9.9) {%
      \textbf{10. \texttt{LoggingInterceptor.preHandle()}}
      — regista o \textit{timestamp} de início em \texttt{request.getAttribute}};
    \node[layer=teal!15]    (ctrl)    at (0,-11.0){%
      \textbf{11. Controlador + Validação Bean} (\texttt{@Valid})
      — desserialização Jackson, restrições JSR-380};
    \node[layer=teal!20]    (svc)     at (0,-12.1){%
      \textbf{12. Serviço de Aplicação} — lógica de domínio, \texttt{Either<AppError,T>}};
    \node[layer=teal!15]    (logpost) at (0,-13.2){%
      \textbf{13. \texttt{LoggingInterceptor.afterCompletion()}}
      — regista método, caminho, estado HTTP e duração em ms};
    \node[layer=orange!10]  (advice)  at (0,-14.3){%
      \textbf{14. \texttt{@RestControllerAdvice}}
      — mapeia exceções não capturadas para respostas de erro estruturadas};

    % Arrows between layers
    \foreach \y/\ny in {0/-1.1,-1.1/-2.2,-2.2/-3.3,-3.3/-4.4,-4.4/-5.5,
                        -5.5/-6.6,-6.6/-7.7,-7.7/-8.8,-8.8/-9.9,-9.9/-11.0,
                        -11.0/-12.1,-12.1/-13.2,-13.2/-14.3} {
      \draw[arrow] (5.6,\y-0.4) -- (5.6,\ny+0.4);
    }
  \end{tikzpicture}
  \caption{\textit{Pipeline} de pedido-resposta do servidor Spring Boot do \textit{Play4Change}}
  \label{fig:pipeline_spring}
\end{figure}

\subsection*{Camada 1 — Nginx}

O Nginx atua como \textit{proxy} inverso à frente do Tomcat.
As suas responsabilidades no \textit{Play4Change} são:
(i)~terminação TLS (o Tomcat vê tráfego HTTP em texto simples internamente);
(ii)~definição dos cabeçalhos de segurança \texttt{X-Frame-Options},
\texttt{X-Content-Type-Options} e \texttt{Referrer-Policy};
(iii)~configuração de \texttt{proxy\_buffering off} e
\texttt{X-Accel-Buffering: no} nos \textit{endpoints} SSE, garantindo
que os eventos são enviados ao cliente sem acumulação no \textit{buffer}
do \textit{proxy}.

\subsection*{Camada 2 — Apache Tomcat (\textit{embedded})}

O Tomcat aceita a ligação TCP, lê os bytes HTTP e entrega um objeto
\texttt{HttpServletRequest} à cadeia de filtros.
O conjunto de fios de execução de pedidos do Tomcat é independente do
\texttt{generationExecutor} da IA: pedidos de longa duração de geração
de tópicos não bloqueiam a capacidade do servidor em processar outros
pedidos.

\subsection*{Camada 3 — \texttt{RateLimitFilter} (\texttt{HIGHEST\_PRECEDENCE})}

Este filtro é o \emph{primeiro componente Spring} a executar
(\texttt{@Order(Ordered.HIGHEST\_PRECEDENCE)} = \texttt{Integer.MIN\_VALUE}).
Aplica-se exclusivamente a caminhos que comecem por \texttt{/auth/}.
O \texttt{RateLimitService} mantém um conjunto de \textit{buckets}
de fichas por IP de cliente (extraído via \texttt{IpExtractor.extractClientIp()},
que respeita o cabeçalho \texttt{X-Forwarded-For} definido pelo Nginx).
Se a ficha não puder ser consumida, a resposta \texttt{429 Too Many
Requests} é emitida com o cabeçalho \texttt{Retry-After} e nenhum
filtro posterior é invocado.

\begin{verbatim}
HTTP/1.1 429 Too Many Requests
Retry-After: 47
Content-Type: application/json

{ "message": "Too many requests. Please wait before trying again." }
\end{verbatim}

\subsection*{Camada 4 — \texttt{RequestIdFilter} (\texttt{@Order(1)})}

Este filtro gera um identificador único por pedido de oito carateres
(extraído de um UUID aleatório) e injeta-o no Contexto de Diagnóstico
Mapeado (MDC) do SLF4J com a chave \texttt{req-id}.
Como o MDC é propagado automaticamente para todas as afirmações de registo
produzidas durante a execução do pedido, qualquer linha de registo do
servidor pode ser correlacionada com um pedido específico sem adicionar
contexto manual.
O identificador é também colocado no cabeçalho de resposta
\texttt{X-Request-Id}, permitindo que os clientes incluam o valor nas
suas próprias comunicações de depuração.
O bloco \texttt{try/finally} garante que o MDC é sempre limpo — omitir
esta limpeza provocaria fugas de valores desatualizados entre pedidos
reutilizando o mesmo fio de execução do conjunto.

\subsection*{Camada 5 — \texttt{CorsFilter} (Spring Security)}

O \texttt{WebMvcConfig} define uma fonte de configuração CORS:
origens permitidas configúráveis via \texttt{app.cors.allowed-origins}
(por omissão: \texttt{localhost:5173} e \texttt{localhost:3000}),
métodos explicitamente listados (\texttt{GET}, \texttt{POST}, \texttt{PUT},
\texttt{PATCH}, \texttt{DELETE}, \texttt{OPTIONS}), e
\texttt{allowCredentials = true} necessário para que o \textit{browser}
envie o \textit{cookie} \texttt{SameSite=Strict} do \textit{token} de
atualização em pedidos entre origens.
Os pedidos \textit{pre-flight} \texttt{OPTIONS} são intercetados aqui e
respondidos diretamente com cabeçalhos \texttt{Access-Control-*} sem
alcançar o controlador.

\subsection*{Camada 6 — \texttt{JwtAuthFilter}}

Este filtro, que estende \texttt{OncePerRequestFilter}, é inserido
\emph{antes} do \texttt{UsernamePasswordAuthenticationFilter} pelo
\texttt{SecurityConfig}:

\begin{verbatim}
.addFilterBefore(jwtAuthFilter,
    UsernamePasswordAuthenticationFilter::class.java)
\end{verbatim}

A lógica é a seguinte:

\begin{enumerate}
  \item Se o caminho pertencer à lista \texttt{ALWAYS\_PUBLIC} ou à lista
    Swagger (em perfil não-produção), o filtro é ignorado e a cadeia
    prossegue.
  \item Se o cabeçalho \texttt{Authorization} começar por \texttt{Bearer },
    o \textit{token} é passado ao \texttt{TokenService.parseAccessToken()}.
    Em caso de sucesso, um \texttt{UsernamePasswordAuthenticationToken}
    com o papel do utilizador é colocado no \texttt{SecurityContextHolder}.
    Qualquer exceção durante a análise (expiração, assinatura inválida,
    \textit{claims} em falta) desencadeia a emissão imediata de um
    \texttt{401} e o retorno — a cadeia de filtros \emph{não} prossegue.
  \item Se não existir cabeçalho \texttt{Authorization}, o contexto é
    deixado vazio e a cadeia continua.
    O Spring Security trata este caso na camada de autorização (Camada 7).
\end{enumerate}

\subsection*{Camada 7 — Autorização Spring Security}

O \texttt{SecurityConfig} define as regras de autorização por rota:

\begin{verbatim}
/auth/**             → permitAll
/api/stats/public    → permitAll
/actuator/health     → permitAll
/actuator/prometheus → permitAll
/swagger-ui/**       → permitAll (não-prod) | hasRole("ADMIN") (prod)
/admin/**            → hasRole("ADMIN")
qualquer outro       → authenticated()
\end{verbatim}

Se o principal autenticado não satisfizer a regra aplicável, o
\texttt{AuthenticationEntryPoint} (sem autenticação → \texttt{401}) ou o
\texttt{AccessDeniedHandler} (autenticado mas papel insuficiente → \texttt{403})
emite a resposta de erro, ambos produzindo um corpo JSON estruturado.

\subsection*{Camada 8 — \texttt{WebMvcMetricsFilter}}

Auto-configurado pelo arranque do Spring Boot Actuator (via dependência
Micrometer), este filtro regista automaticamente o contador e o histograma
\texttt{http\_server\_requests\_seconds} com as etiquetas \texttt{method},
\texttt{uri} (padrão de rota, e.g. \texttt{/admin/topics/\{id\}}),
\texttt{status} e \texttt{outcome}.
O \texttt{LoggingInterceptor} \emph{não} duplica este registo de propósito —
os comentários do código explicam que dois registos com conjuntos de
chaves de etiquetas diferentes corromperiam a família de métricas no
Prometheus.

\subsection*{Camada 9 — \texttt{DispatcherServlet}}

O \texttt{DispatcherServlet} realiza o mapeamento de controladores
(resolvendo o método anotado com \texttt{@PostMapping} / \texttt{@GetMapping}
que corresponde ao padrão URI e ao método HTTP), a negociação de conteúdo
(seleccionando Jackson para corpos JSON, \texttt{SseEmitter} para fluxos
SSE) e a resolução de argumentos de método (parâmetros \texttt{@PathVariable},
\texttt{@RequestBody}, \texttt{@AuthenticationPrincipal}).

\subsection*{Camada 10 — \texttt{LoggingInterceptor.preHandle()}}

O Spring MVC invoca \texttt{preHandle()} em cada interceptor registado
\emph{antes} de o método do controlador ser executado.
O \texttt{LoggingInterceptor} regista o \textit{timestamp} de início de
milissegundos em \texttt{request.setAttribute("request.timer.start", now)}
e emite uma linha de registo de nível \texttt{DEBUG}:
\texttt{"--> POST /admin/topics"}.

\subsection*{Camada 11 — Controlador + Validação Bean}

O método do controlador é invocado com os argumentos já resolvidos.
Quando um parâmetro está anotado com \texttt{@Valid}, o Hibernate Validator
valida todas as restrições JSR-380 (e.g., \texttt{@NotBlank},
\texttt{@Min(1)}) antes de invocar o corpo do método.
Uma falha de validação lança \texttt{MethodArgumentNotValidException}, que
é capturada pelo \texttt{@RestControllerAdvice} (Camada 14) e mapeada para
uma resposta \texttt{400 Bad Request} com a lista de erros de campo.

\subsection*{Camada 12 — Serviço de Aplicação}

O controlador delega no serviço de aplicação, que executa a lógica de
domínio devolvendo \texttt{Either<AppError, T>}.
O controlador desdobra o resultado: o ramo \texttt{Right} é mapeado para
uma resposta de sucesso; o ramo \texttt{Left} é mapeado para o código de
estado HTTP definido pelo subtipo de \texttt{AppError}
(ver Tabela~\ref{tab:codigos_erro}).
Esta abordagem torna os percursos de erro \emph{estruturalmente visíveis}
em tempo de compilação — sem tratamento de exceções implícito que possa
ser silenciosamente ignorado.

\subsection*{Camada 13 — \texttt{LoggingInterceptor.afterCompletion()}}

Após o método do controlador completar (com sucesso ou com exceção),
\texttt{afterCompletion()} calcula a duração decorrida, recuperando o
\textit{timestamp} armazenado pelo \texttt{preHandle()}, e emite a linha
de registo final:
\texttt{"<-- POST /admin/topics → 201 [143ms]"}.
Esta linha, combinada com o \texttt{req-id} no MDC, fornece uma rastreabilidade
completa de ponta a ponta de cada pedido nos registos estruturados.

\subsection*{Camada 14 — \texttt{@RestControllerAdvice}}

O \texttt{AuthExceptionHandler}, anotado com \texttt{@RestControllerAdvice},
trata as exceções não capturadas que escapam ao controlador.
Os mapeamentos cobrem:

\begin{itemize}
  \item \texttt{IllegalArgumentException} → \texttt{400}
    (erros de validação de domínio não cobertos por Bean Validation)
  \item \texttt{SecurityException} → \texttt{401}
    (violações de segurança explícitas)
  \item \texttt{IllegalStateException} → \texttt{502 Bad Gateway}
    (erros de serviços \textit{upstream} — o corpo \emph{não} expõe a
    mensagem de exceção original para evitar fugas de URLs ou detalhes
    internos de fornecedores)
  \item \texttt{HttpMessageNotReadableException} → \texttt{400}
    (corpo de pedido JSON malformado)
  \item \texttt{MethodArgumentNotValidException} → \texttt{400}
    (falha de Bean Validation, com lista detalhada de erros de campo)
  \item \texttt{ObjectOptimisticLockingFailureException} → \texttt{409}
    (colisão de bloqueio otimista em operações concorrentes)
\end{itemize}

\subsection*{Fluxo paralelo: executor de geração assíncrona}

Os pedidos de criação de tópico e de processamento de dificuldade não
bloqueiam o fio de execução do Tomcat para aguardar o resultado da IA.
O \texttt{AsyncConfig} configura um \texttt{ThreadPoolTaskExecutor}
dedicado com o nome \texttt{generationExecutor}:

\begin{verbatim}
@Bean("generationExecutor")
fun generationExecutor(): Executor = ThreadPoolTaskExecutor().apply {
    corePoolSize  = poolSize      // padrão: 3
    maxPoolSize   = poolSize      // sem crescimento adicional
    setQueueCapacity(25)          // máximo de 25 pedidos em espera
    setThreadNamePrefix("gen-")   // visível em JVM thread dumps
    initialize()
}
\end{verbatim}

Quando a fila atingir a capacidade (25 tarefas pendentes), novos pedidos
de geração são rejeitados com \texttt{503 Service Unavailable} e o alerta
Prometheus \texttt{GenerationThreadPoolQueueSaturated} é disparado.
Os nomes dos fios de execução \texttt{gen-1}, \texttt{gen-2}, \texttt{gen-3}
são visíveis em despejos de pilha da JVM, facilitando o diagnóstico de
bloqueios ou fugas de fios de execução.

\begin{figure}[htbp]
  \centering
  \begin{tikzpicture}[>=Stealth, font=\small,
    comp/.style={draw, rounded corners=3pt, fill=#1,
                 minimum width=3.6cm, minimum height=0.9cm, align=center},
    note/.style={font=\tiny, align=center}
  ]
    % Synchronous path
    \node[comp=orange!15] (req)   at (0, 0)    {Pedido HTTP};
    \node[comp=blue!10]   (filt)  at (0,-1.5)  {Filtros (1--8)};
    \node[comp=green!10]  (ctrl)  at (0,-3.0)  {Controlador};
    \node[comp=teal!10]   (svc)   at (0,-4.5)  {Serviço de Aplicação};
    \node[comp=orange!10] (resp)  at (0,-6.0)  {\texttt{201 Created}};

    % Async path
    \node[comp=yellow!20] (exec)  at (5.5,-3.0) {\texttt{generateAsync()}\\submete tarefa};
    \node[comp=purple!10] (pool)  at (5.5,-4.5) {\texttt{gen-1..3}\\conjunto de fios};
    \node[comp=red!10]    (llm)   at (5.5,-6.0) {LLM Mistral\\(IA)};
    \node[comp=blue!10]   (sse)   at (5.5,-7.5) {Eventos SSE\\para clientes};

    % Sync flow
    \draw[->] (req)  -- (filt);
    \draw[->] (filt) -- (ctrl);
    \draw[->] (ctrl) -- (svc);
    \draw[->] (svc)  -- (resp);

    % Branch to async
    \draw[->, dashed] (svc) -- node[above,font=\tiny]{async} (exec);
    \draw[->] (exec) -- (pool);
    \draw[->] (pool) -- (llm);
    \draw[->] (llm)  -- (sse);

    % Labels
    \node[note] at (-2.2,-3.0) {fio Tomcat\\(bloqueante)};
    \node[note] at (7.8,-4.5) {fio \texttt{gen-*}\\(independente)};
  \end{tikzpicture}
  \caption{Separação entre o percurso síncrono HTTP e o executor de geração assíncrona}
  \label{fig:sinc_async}
\end{figure}

Esta separação é a razão pela qual a criação de um tópico devolve
imediatamente \texttt{201 Created}: o fio de execução do Tomcat retorna
ao conjunto assim que o trabalho de geração é submetido ao
\texttt{generationExecutor}, estando disponível para servir outros pedidos
enquanto os fios \texttt{gen-*} se ocupam da chamada ao LLM.

% ─────────────────────────────────────────────────────────────────────────────
\section{Portal de Administração: Criação de Tópico a partir de URL}
\label{sec:topic_criacao_url}

A criação de um tópico a partir de um endereço URL constitui o fluxo
primário de autoria de conteúdo para os administradores. Envolve o
\textit{frontend} web, a API REST, dois serviços externos (extrator de
conteúdo e Mistral AI), o armazenamento de objetos MinIO e o PostgreSQL.

\subsection*{Pedido HTTP}

\begin{verbatim}
POST /admin/topics
Authorization: Bearer <token_de_acesso>
Content-Type: application/json

{
  "title":       "Introdução à Mobilidade Urbana",
  "description": "Conceitos fundamentais de transporte sustentável",
  "category":    "Cidade Inteligente",
  "url":         "https://example.org/guia-mobilidade",
  "taskCount":   15,
  "difficulty":  "BEGINNER",
  "language":    "pt-PT"
}
\end{verbatim}

\subsection*{Validação}

O \texttt{TopicController.createFromUrl()} valida o corpo do pedido via
Bean Validation (\texttt{@Valid}).
O \texttt{TopicManagementService} executa então a validação de domínio:
\texttt{taskCount > 0} (caso contrário, \texttt{400 Bad Request} com
\texttt{\{field:"taskCount", reason:"must be greater than 0"\}}).

\subsection*{Extração de conteúdo}

O \texttt{ContentExtractorPort.extractFromUrl(url)} obtém a página, remove
as marcações HTML e devolve texto simples.
Se a obtenção falhar — erro de rede, 404, corpo vazio — o serviço levanta
\texttt{BadRequest.InvalidField("url", "could not fetch content: \ldots")}
que o controlador mapeia para uma resposta \texttt{400}.
Se o texto extraído for vazio (página renderizada pelo lado do cliente sem
conteúdo no servidor), é devolvido um \texttt{400} distinto:
\texttt{"url returned no extractable text"}.

\subsection*{Armazenamento de objetos}

O texto extraído é carregado para o MinIO com a chave
\texttt{topics/\{topicId\}/content.txt} via \texttt{FileStoragePort}.
Uma falha no MinIO levanta \texttt{InternalServerError.UnexpectedException}
→ \texttt{500}.
O URL do objeto MinIO devolvido é guardado como \texttt{Topic.contentSourceRef}.

\subsection*{Registo na base de dados}

É guardada uma entidade \texttt{Topic} com:

\begin{verbatim}
id            = <UUID>
status        = PENDING
currentPhase  = INGESTION
createdBy     = <adminId do JWT>
language      = "pt-PT"
audienceLevel = BEGINNER
taskCount     = 15
\end{verbatim}

\subsection*{Resposta imediata}

\begin{verbatim}
HTTP/1.1 201 Created
{
  "id": "<topicId>",
  "status": "PENDING",
  "currentPhase": "INGESTION",
  ...
}
\end{verbatim}

O identificador do tópico é devolvido \emph{imediatamente} — o
\textit{pipeline} de geração por IA é executado de forma \emph{assíncrona},
pelo que a resposta HTTP não é bloqueada.

\subsection*{\textit{Pipeline} de geração assíncrona}

O \texttt{orchestrator.generateAsync(topicId)} submete uma tarefa ao
executor \texttt{@Async("generationExecutor")}, um conjunto de
\textit{threads} delimitado (capacidade da fila: 25) dedicado ao trabalho
de geração.
O orquestrador conduz o tópico através de quatro fases
(Figura~\ref{fig:pipeline_geracao}).

\begin{figure}[htbp]
  \centering
  \begin{tikzpicture}[>=Stealth, font=\small,
    phase/.style={draw, rounded corners=4pt, fill=yellow!20,
                  minimum width=2.4cm, minimum height=1cm, align=center},
    arrow/.style={->, thick}
  ]
    \node[phase] (ing)  at (0,0)    {INGESTION\\{\tiny obter e guardar}};
    \node[phase] (ana)  at (3.3,0)  {ANALYSIS\\{\tiny extrair objetivos}};
    \node[phase] (gen)  at (6.6,0)  {GENERATION\\{\tiny LLM + dedup}};
    \node[phase] (idx)  at (9.9,0)  {INDEXING\\{\tiny embeder e guardar}};
    \node[phase,fill=green!20]  (act)  at (6.6,-2.2) {ACTIVE};
    \node[phase,fill=red!20]    (fail) at (3.3,-2.2) {FAILED};

    \draw[arrow] (ing) -- (ana);
    \draw[arrow] (ana) -- (gen);
    \draw[arrow] (gen) -- (idx);
    \draw[arrow] (idx) -- (act);
    \draw[arrow, dashed, color=red] (gen) -- node[right,font=\tiny]{erro IA} (fail);
    \draw[arrow, dashed, color=red] (ana) -- node[left,font=\tiny]{\textit{timeout}} (fail);
  \end{tikzpicture}
  \caption{Fases do \textit{pipeline} de geração de tópicos}
  \label{fig:pipeline_geracao}
\end{figure}

\paragraph{INGESTION → ANALYSIS.}
O \texttt{PhaseTransitionService.transitionTo(ANALYSIS)} regista a duração
gasta na fase INGESTION no \texttt{TopicPhaseLog}, atualiza
\texttt{Topic.currentPhase} e publica um evento SSE \texttt{phase-change}
para os clientes a escutar.

\paragraph{ANALYSIS → GENERATION.}
Uma passagem de análise extrai o objetivo do módulo a partir do texto,
utilizado posteriormente no \textit{prompt} de IA.

\paragraph{GENERATION.}
A chamada central ao LLM.
O \texttt{LangChain4jTaskGenerationAdapter.generateTasks()} monta duas
mensagens (ver Secção~\ref{sec:engenharia_prompts}) e invoca
\texttt{chatModel.generate()}.
O adaptador analisa o \textit{array} JSON de resposta, descartando os
elementos que falham a validação.
Se houver menos tarefas válidas do que o solicitado, \textbf{é realizada
uma única nova tentativa} com o \textit{prompt} original acrescido de um
parágrafo \texttt{schemaReminder}:

\begin{verbatim}
Your previous response was incomplete or malformed.
Please respond with EXACTLY 15 tasks in the JSON array
schema above. Every task must have: title, description,
hint, pointsReward, options (4 items), and
correctAnswerIndex set to 0. No extra fields. Valid JSON only.
\end{verbatim}

Se, após a nova tentativa, não existir nenhuma tarefa válida, o adaptador
devolve
\texttt{Either.Left(ServiceUnavailable.DependencyUnavailable("mistral-ai-schema"))}
e o orquestrador transita o tópico para \texttt{FAILED}.

Cada tarefa analizada com sucesso é verificada quanto a duplicação
semântica via
\texttt{PgVectorDeduplicationService.isDuplicate()} (similaridade do
cosseno $> 0{,}90$ no espaço de \textit{embeddings} do mesmo módulo).
Tarefas duplicadas são marcadas \texttt{GenerationStatus.DUPLICATE} e
excluídas da persistência; as tarefas inéditas são guardadas como
registos \texttt{TaskTemplate} e emitidas para o fluxo SSE como eventos
\texttt{generation-progress}.

\paragraph{INDEXING.}
Os \textit{embeddings} das tarefas são persistidos na coluna
\textit{pgvector} de \texttt{task\_templates} via \texttt{JdbcTemplate}
(o tipo da extensão vectorial requer o cast \texttt{::vector} no SQL).

\paragraph{ACTIVE.}
O \texttt{Topic.status} é atualizado para \texttt{ACTIVE} e um evento
SSE final \texttt{complete} é emitido.

% ─────────────────────────────────────────────────────────────────────────────
\section{Portal de Administração: Criação de Tópico a partir de PDF}
\label{sec:topic_criacao_pdf}

O percurso de PDF difere do percurso de URL apenas na extração de conteúdo
e na codificação HTTP.

\begin{verbatim}
POST /admin/topics/pdf
Authorization: Bearer <token_de_acesso>
Content-Type: multipart/form-data

--boundary
Content-Disposition: form-data; name="title"
Introdução à Mobilidade Urbana

--boundary
Content-Disposition: form-data; name="file"; filename="guia.pdf"
Content-Type: application/pdf
<bytes binários do PDF>
\end{verbatim}

O controlador delega no \texttt{TopicManagementService.createFromPdf()}.
O \texttt{ContentExtractorPort.extractFromPdf(bytes)} utiliza o Apache
PDFBox para extrair o texto do binário; um resultado vazio (PDF de imagem
digitalizada sem OCR) devolve \texttt{400 "PDF contains no extractable
text"}.
Em caso de sucesso, os bytes do PDF são carregados para o MinIO com a
chave \texttt{topics/\{topicId\}/content.pdf} (MIME:
\texttt{application/pdf}) e o \textit{pipeline} prossegue de forma
idêntica ao percurso de URL.

% ─────────────────────────────────────────────────────────────────────────────
\section{Observação da Geração: \textit{Server-Sent Events}}
\label{sec:sse_deepdive}

Após a resposta \texttt{201}, o \textit{frontend} web armazena o
\texttt{topicId} devolvido e ativa o \textit{hook}
\texttt{useTopicProgress(topicId)}, que abre uma ligação SSE:

\begin{verbatim}
GET /admin/topics/{id}/progress
Authorization: Bearer <token_de_acesso>
Accept: text/event-stream
Cache-Control: no-cache
\end{verbatim}

O servidor responde com \texttt{Content-Type: text/event-stream} e
\texttt{X-Accel-Buffering: no} — instrução ao \textit{proxy} inverso
Nginx para desativar o \textit{buffering} de resposta, garantindo que os
eventos são enviados imediatamente em vez de acumulados.

O \textit{hook} cliente lê o fluxo linha a linha usando a API
\texttt{ReadableStream} (sem biblioteca \textit{EventSource} de terceiros),
analisando o formato de fio SSE padrão do W3C~\cite{w3c2021sse}:

\begin{verbatim}
event: phase-change
data: {"phase":"ANALYSIS"}

event: generation-progress
data: {"completed":7,"total":15}

event: complete
data: {}
\end{verbatim}

O componente React \texttt{TopicProgressStepper} volta a renderizar a cada
alteração de estado, mostrando a fase atual e uma barra de progresso.
Se a ligação SSE cair (erro de rede, reinício do servidor), o \textit{hook}
define \texttt{failed: true} — salvo se o \texttt{AbortError} tiver origem
no efeito de limpeza — e a navegação é redirecionada para a página de
detalhe do tópico, onde está disponível a opção de nova tentativa.

O \texttt{SseTopicEventPublisher} do lado do servidor mantém um
\texttt{Map<String, SseEmitter>}; um \textit{emitter} é registado em cada
chamada \texttt{GET /progress} e removido quando o \textit{emitter}
termina ou expira por \textit{timeout}.
Se o tópico já se encontrar \texttt{ACTIVE} quando a ligação SSE chegar,
\texttt{ssePublisher.completed()} dispara imediatamente; se já estiver
\texttt{FAILED}, \texttt{ssePublisher.failed()} dispara com a razão de
falha armazenada.

% ─────────────────────────────────────────────────────────────────────────────
\section{Portal de Administração: Gestão de Utilizadores}
\label{sec:gestao_utilizadores}

\begin{table}[ht]
  \caption{Pontos de extremidade de gestão de utilizadores do administrador}
  \label{tab:admin_user_endpoints}
  \vskip0.5em
  \centering
  \begin{tabular}{llp{5.8cm}}
    \toprule
    \textbf{Método e Caminho} & \textbf{Estado} & \textbf{Resposta / Comportamento} \\
    \midrule
    \texttt{GET /admin/users?page=\&size=} & 200 &
      \texttt{AdminUserPage} (lista paginada com e-mail, papel, data de registo) \\
    \texttt{GET /admin/users/\{id\}} & 200 / 404 &
      \texttt{AdminUserDetail} (perfil e estatísticas) \\
    \texttt{POST /admin/users/\{id\}/promote} & 200 / 404 &
      Eleva o papel do utilizador para \texttt{ADMIN}; devolve
      \texttt{AdminUserFull} atualizado \\
    \texttt{GET /admin/users/\{id\}/enrollments} & 200 &
      \texttt{AdminUserEnrollment[]} com estado, título do tópico e
      percentagem de progresso \\
    \texttt{GET /admin/users/\{id\}/badges} & 200 &
      \texttt{AdminUserBadge[]} com nome do \textit{badge}, data de
      conquista e tópico associado \\
    \texttt{GET /admin/users/\{id\}/enrollments/\{eid\}/roadmap} & 200 &
      \texttt{AdminRoadmapNode[]} — tarefas ordenadas com estado de
      conclusão \\
    \bottomrule
  \end{tabular}
\end{table}

A operação de \textbf{promoção} é a mais sensível: altera o âmbito de
permissões do utilizador para todos os pedidos subsequentes.
O Spring Security lê o papel a partir das \textit{claims} do JWT; o novo
papel produz efeito no \emph{próximo} início de sessão (o \textit{token}
de acesso atual reflete o papel no momento da emissão).

% ─────────────────────────────────────────────────────────────────────────────
\section{Portal de Administração: Gestão de Tarefas e Grafo de Aprendizagem}
\label{sec:tarefas_grafo}

\subsection*{Visualização e edição de tarefas}

\begin{verbatim}
GET  /admin/topics/{id}/tasks               → TaskTemplate[]
GET  /admin/topics/{id}/struggle-tasks      → AdaptiveTaskAdmin[]
GET  /admin/topics/{id}/struggle-path-stats → StrugglePathStats[]
PUT  /admin/tasks/{templateId}              → TaskTemplate (atualizado)
PUT  /admin/adaptive-tasks/{taskId}         → AdaptiveTaskAdmin (atualizado)
\end{verbatim}

Um \texttt{PUT} a um modelo de tarefa incrementa
\texttt{TaskTemplate.version}, valor que é transportado para os novos
registos de \texttt{TaskAssignment}.
Os aprendentes que já possuem uma atribuição sobre a versão anterior não
são afetados; apenas as novas atribuições criadas após a atualização
apresentam o conteúdo editado.

\subsection*{Pré-requisitos e deteção de ciclos}

\begin{verbatim}
GET  /admin/topics/{id}/prerequisites
POST /admin/topics/{id}/prerequisites
     { "prerequisiteIds": ["<id1>","<id2>"] }
GET  /admin/learning-graph  → { edges: [{topicId, prerequisiteId}] }
\end{verbatim}

O \texttt{TopicManagementService.setPrerequisites()} executa um algoritmo
de deteção de ciclos baseado em DFS antes de persistir.
Se a adição dos pré-requisitos solicitados criasse um ciclo no grafo de
aprendizagem, é devolvido:

\begin{verbatim}
400 Bad Request
{
  "field": "prerequisiteIds",
  "reason": "setting these prerequisites would create a cycle
             in the learning graph"
}
\end{verbatim}

O \textit{frontend} web renderiza o grafo completo com
\texttt{LearningPathDagView}, uma visualização DAG personalizada em SVG.

% ─────────────────────────────────────────────────────────────────────────────
\section{Jornada do Aprendente: Inscrição num Tópico}
\label{sec:inscricao_deepdive}

A inscrição é iniciada pelo aprendente na aplicação móvel através do ecrã
\textbf{Explore} (\texttt{ExploreScreen}), que lista os tópicos ativos
disponíveis.
O \texttt{HttpExploreRepository} emite o seguinte pedido quando o utilizador
toca em «Inscrever»:

\begin{verbatim}
POST /user/enroll
Authorization: Bearer <token_do_aprendente>
Content-Type: application/json

{ "topicId": "<uuid>" }
\end{verbatim}

O \texttt{EnrollmentService.enroll()} aplica quatro guardas em sequência:

\begin{enumerate}
  \item \textbf{O tópico tem de estar ACTIVE.}
    Se o estado for \texttt{PENDING}, \texttt{GENERATING} ou \texttt{FAILED}:
    \texttt{400 InvalidField("topicId","topic is not active")}.

  \item \textbf{O tópico não pode ter expirado.}
    \texttt{Topic.isExpired()} verifica \texttt{expiresAt < now}:
    \texttt{400 InvalidField("topicId","topic has expired")}.

  \item \textbf{Todos os pré-requisitos têm de estar concluídos.}
    Se algum pré-requisito não constar no conjunto de inscrições concluídas
    do utilizador:
    \texttt{400 InvalidField("topicId","prerequisites not completed: [\ldots]")}.

  \item \textbf{Sem inscrição duplicada.}
    Se já existir uma inscrição ativa: \texttt{409 Conflict}.
    Se existir uma inscrição \emph{em pausa}, o serviço reativa-a em vez
    de criar uma duplicada.
\end{enumerate}

Em caso de sucesso, o \texttt{EnrollmentService} cria o registo de
\texttt{Enrollment} e imediatamente a \emph{primeira atribuição de tarefa}
para \texttt{dayIndex = 0}.
A ordem das opções é misturada deterministicamente usando
\texttt{TaskShuffleSeed.shuffleOptions(optionCount, userId, taskId,
enrollmentId)}, garantindo que o mesmo utilizador vê sempre a mesma
mistura ao voltar à plataforma enquanto utilizadores distintos veem
ordenações diferentes — eliminando o reconhecimento trivial de padrões
que a ordem fixa de opções possibilita.
O contador Prometheus \texttt{topic\_enrollments\_total} é incrementado.

\begin{verbatim}
HTTP/1.1 201 Created
{
  "id": "<enrollmentId>",
  "topicId": "<topicId>",
  "status": "ACTIVE",
  "currentDayIndex": 0,
  "streakDays": 0,
  ...
}
\end{verbatim}

% ─────────────────────────────────────────────────────────────────────────────
\section{Jornada do Aprendente: Obtenção e Submissão de Tarefa (App Móvel)}
\label{sec:submissao_dificuldade}

\subsection*{Obtenção da tarefa diária}

O ecrã \texttt{Home} apresenta ao aprendente o tópico em que está inscrito.
Quando o utilizador toca em «Resolver tarefa de hoje», o componente dispara
\texttt{HomeEffect.NavigateToTask(userTaskId)} e o
\texttt{DefaultRootComponent} empurra \texttt{Config.Task(userTaskId)}
para a pilha de navegação.
O \texttt{HttpTaskRepository.getTask(topicId)} emite:

\begin{verbatim}
GET /tasks/today?topicId={topicId}
Authorization: Bearer <token_do_aprendente>
X-Timezone: Europe/Lisbon
\end{verbatim}

O cabeçalho \texttt{X-Timezone} é definido com
\texttt{TimeZone.currentSystemDefault().id} do Kotlin Multiplatform, que
resolve o fuso horário local sem depender de APIs de plataforma adicionais.
O servidor usa-o para calcular se a atribuição do dia corrente já foi
gerada (as atribuições são geradas por dia de calendário no fuso horário
do aprendente).

O servidor pode responder com três estados especiais além de
\texttt{200 OK}:

\begin{table}[ht]
  \caption{Respostas especiais de \texttt{GET /tasks/today} tratadas pelo cliente móvel}
  \label{tab:task_estados}
  \vskip0.5em
  \centering
  \begin{tabular}{llp{6.5cm}}
    \toprule
    \textbf{Estado HTTP} & \textbf{\texttt{NetworkError}} & \textbf{Comportamento no ecrã} \\
    \midrule
    \texttt{202 Accepted} & \texttt{TaskGenerationPending} &
      Mostra estado de espera «O conteúdo ainda está a ser gerado» \\
    \texttt{404 Not Found} & \texttt{NoTaskAvailable} &
      Mostra estado vazio «Já concluíste a tarefa de hoje!» \\
    \texttt{409 Conflict} & \texttt{StruggleOpen(enrollmentId)} &
      Emite \texttt{TaskEffect.NavigateToStruggle(enrollmentId)}\\
    \bottomrule
  \end{tabular}
\end{table}

O cabeçalho \texttt{X-Open-Struggle-Enrollment} da resposta \texttt{409}
transporta o \texttt{enrollmentId} necessário para abrir a sessão de
remediação ativa.

\subsection*{Submissão da resposta}

O aprendente seleciona uma das quatro opções apresentadas.
O \texttt{HttpTaskRepository.submitAnswer()} emite:

\begin{verbatim}
POST /user/tasks/{assignmentId}/submit
Authorization: Bearer <token_do_aprendente>
Content-Type: application/json
X-Timezone: Europe/Lisbon

{ "selectedOption": 2 }
\end{verbatim}

O \texttt{TaskService} resolve a atribuição, consulta a ordem misturada
de opções para determinar a opção canónica selecionada, compara com
\texttt{TaskTemplate.correctAnswer} e atualiza a atribuição:

\begin{verbatim}
submittedAt       = agora
status            = COMPLETED
isCorrect         = true | false
pointsAwarded     = template.pointsReward | 0
wrongAttemptCount += 1   (se incorreto)
\end{verbatim}

\paragraph{Deteção de dificuldade.}
Se \texttt{wrongAttemptCount $\geq$ 2} e o aprendente não tiver ainda
um ramo adaptativo atribuído para esta tarefa,
\texttt{HandleStruggleService.triggerAsync()} é invocado.
O padrão de erro é classificado com base na opção selecionada:

\begin{itemize}
  \item \texttt{WRONG\_CONCEPT} — resposta consistentemente errada, sem
    proximidade à correta
  \item \texttt{PARTIAL\_UNDERSTANDING} — opção próxima escolhida
    repetidamente
  \item \texttt{READING\_ERROR} — conceito correto mas formulação errada
  \item \texttt{TIME\_PRESSURE} — submissão rápida sem padrão consistente
\end{itemize}

A resposta de submissão ao aprendente é \texttt{200 OK} com o resultado;
a geração adaptativa decorre em segundo plano sem bloquear a resposta.

\paragraph{Tratamento da dificuldade no cliente móvel.}
O DTO de submissão inclui o campo \texttt{struggleTriggered: Boolean}.
O \texttt{DefaultTaskComponent} verifica este campo: se for \texttt{true},
emite o efeito \texttt{TaskEffect.NavigateToStruggle(enrollmentId)}.
O \texttt{DefaultRootComponent} trata este efeito empurrando
\texttt{Config.Struggle(enrollmentId)} para a pilha, levando o aprendente
ao ecrã de remediação adaptativa (Figura~\ref{fig:mobile_task_flow}).

\begin{figure}[htbp]
  \centering
  \begin{tikzpicture}[>=Stealth, font=\small,
    scr/.style={draw, fill=#1, rounded corners=4pt,
                minimum width=3.2cm, minimum height=0.9cm, align=center},
    ev/.style={font=\tiny, align=center, text width=2.8cm}
  ]
    \node[scr=green!12]   (home)     at (0, 0)   {Home\\{\tiny lista de tópicos}};
    \node[scr=blue!10]    (task)     at (0,-2.2) {Task\\{\tiny 4 opções embaralhadas}};
    \node[scr=orange!15]  (result)   at (-3.5,-4.2){Resultado\\{\tiny correto / errado}};
    \node[scr=red!10]     (struggle) at ( 3.5,-4.2){Struggle\\{\tiny ramo adaptativo}};
    \node[scr=gray!12]    (empty)    at (0,-4.2) {Vazio\\{\tiny já concluído}};

    \draw[->] (home)   -- node[left,ev]{NavigateToTask} (task);
    \draw[->] (task)   -- node[left,ev]{isCorrect=true\\ou errado\\ $<$2 vezes} (result);
    \draw[->] (task)   -- node[right,ev]{struggleTriggered\\=true} (struggle);
    \draw[->] (task.south) -- node[right,ev]{404 NoTask\\Available} (empty);
    \draw[->] (result) to[bend right=35] node[left,ev]{nova tarefa\\amanhã} (home);
    \draw[->] (struggle) to[bend left=35] node[right,ev]{session\\Resolved} (home);
  \end{tikzpicture}
  \caption{Fluxo de navegação na jornada de tarefa do aprendente (app móvel)}
  \label{fig:mobile_task_flow}
\end{figure}

% ─────────────────────────────────────────────────────────────────────────────
\section{Personalização por IA: O \textit{Pipeline} de Ramos Adaptativos}
\label{sec:pipeline_adaptativo}

Quando é detetada uma dificuldade, o sistema inicia a personalização
adaptativa (Figura~\ref{fig:fluxo_adaptativo}).

\begin{figure}[htbp]
  \centering
  \begin{tikzpicture}[>=Stealth, font=\small,
    box/.style={draw, rounded corners=3pt, fill=blue!8,
                minimum width=3.2cm, minimum height=0.8cm, align=center},
    decision/.style={draw, diamond, fill=yellow!20,
                     minimum width=2.8cm, minimum height=0.8cm, align=center},
    result/.style={draw, rounded corners=3pt, fill=green!15,
                   minimum width=2.8cm, minimum height=0.7cm, align=center}
  ]
    \node[box] (embed)  at (0, 0)   {1. Embeder dificuldade\\{\tiny 1024 dim. — Mistral}};
    \node[box] (search) at (0,-1.8) {2. Pesquisa ANN pgvector\\{\tiny distância do cosseno}};
    \node[decision] (sim) at (0,-3.6) {similaridade?};

    \node[result] (full)  at (-4.2,-5.4) {FULL\_REUSE\\{\tiny sim $>$ 0{,}90}};
    \node[result] (part)  at (0,-5.4)    {PARTIAL\_REUSE\\{\tiny 0{,}65 -- 0{,}90}};
    \node[result] (fresh) at (4.2,-5.4)  {FRESH\_GENERATION\\{\tiny $<$ 0{,}65}};

    \node[box] (persist) at (0,-7.2) {3. Persistir ramo e \textit{embedding}};
    \node[box] (save)    at (0,-9.0) {4. Guardar StruggleSession\\com AdaptiveTasks};

    \draw[->] (embed) -- (search);
    \draw[->] (search) -- (sim);
    \draw[->] (sim) -| node[above,font=\tiny]{elevada} (full);
    \draw[->] (sim) -- node[right,font=\tiny]{média} (part);
    \draw[->] (sim) -| node[above,font=\tiny]{baixa} (fresh);
    \draw[->] (full)  |- (persist);
    \draw[->] (part)  -- (persist);
    \draw[->] (fresh) |- (persist);
    \draw[->] (persist) -- (save);
  \end{tikzpicture}
  \caption{Fluxo de decisão na geração de ramos adaptativos}
  \label{fig:fluxo_adaptativo}
\end{figure}

\subsection*{Passo 1 — Incorporação da dificuldade}

É construída uma representação textual da dificuldade:
\texttt{"\{errorPattern\} \{taskDescription\} \{subjectDomain\}"}.
Este texto é incorporado pelo modelo de \textit{embedding} do Mistral,
produzindo um vetor de 1024 dimensões em ponto flutuante.
O \textit{embedding} é validado: se a sua dimensionalidade diferir de 1024
(por exemplo, por alteração de versão da API do modelo), o adaptador lança
uma \texttt{IllegalArgumentException} que é capturada e mapeada para
\texttt{ServiceUnavailable}.

\subsection*{Passo 2 — Pesquisa por similaridade}

O \texttt{PgVectorDeduplicationService.findSimilarStruggle()} executa:

\begin{verbatim}
SELECT ab.id AS branch_id,
       1 - (se.embedding <=> ?::vector) AS similarity
FROM   struggle_events se
JOIN   adaptive_branches ab ON ab.struggle_event_id = se.id
WHERE  ab.status = 'COMPLETED'
ORDER  BY se.embedding <=> ?::vector
LIMIT  1;
\end{verbatim}

O operador \texttt{<=>} é o operador de distância do cosseno do
\textit{pgvector}.
O \texttt{1 -} converte distância em similaridade.
Um índice aproximado de vizinho mais próximo (IVFFlat) na coluna
\texttt{struggle\_events.embedding} mantém esta consulta abaixo do
milissegundo mesmo à escala~\cite{johnson2019billion}.

\subsection*{Passo 3 — Estratégia de reutilização}

\begin{description}
  \item[\texttt{FULL\_REUSE} (similaridade $>$ 0{,}90).]
    O ramo adaptativo armazenado é carregado diretamente a partir de
    \texttt{adaptive\_tasks}.
    Não é realizada nenhuma chamada ao LLM; o sistema explora
    integralmente a geração anterior.
    Este é o \textit{percurso quente} — custo zero de \textit{tokens}.

  \item[\texttt{PARTIAL\_REUSE} (0{,}65 -- 0{,}90).]
    Uma fração do ramo armazenado (proporcional ao valor de similaridade)
    é reutilizada; a restante parte é gerada de novo via
    \texttt{StruggleAnalysisPrompt}.
    Um novo registo de ramo é persistido combinando ambos os conjuntos.

  \item[\texttt{FRESH\_GENERATION} ($<$ 0{,}65).]
    Não existe nenhuma dificuldade suficientemente similar.
    É realizada uma chamada completa ao LLM usando
    \texttt{StruggleAnalysisPrompt}, o resultado é persistido como um
    novo ramo e o seu \textit{embedding} é armazenado de forma a que
    utilizadores futuros com dificuldades similares possam beneficiar.
\end{description}

Cada estratégia de reutilização é rastreada pelo contador Prometheus
\texttt{ai.adaptive.reuse\{strategy="full"|"partial"|"fresh"\}},
permitindo observar a taxa de crescimento do mecanismo de personalização.

% ─────────────────────────────────────────────────────────────────────────────
\section{Jornada do Aprendente: Remediação Adaptativa no Cliente Móvel}
\label{sec:struggle_mobile}

Quando o aprendente chega ao ecrã \textbf{Struggle}, o
\texttt{DefaultStruggleComponent} carrega a sessão de remediação ativa:

\begin{verbatim}
GET /struggle/enrollment/{enrollmentId}
Authorization: Bearer <token_do_aprendente>
\end{verbatim}

O \texttt{HttpStruggleRepository.getSession()} desserializa o
\texttt{StruggleSessionDto}, mapeando cada \texttt{AdaptiveTaskDto} para
um modelo de domínio \texttt{AdaptiveTask}.
A sessão contém o campo \texttt{errorPattern} (e.g.,
\texttt{"WRONG\_CONCEPT"}), que o ecrã usa para apresentar uma mensagem
contextualizada ao aprendente.

O aprendente resolve cada tarefa adaptativa de forma sequencial.
Cada submissão envia:

\begin{verbatim}
POST /struggle/{sessionId}/tasks/{taskId}/submit
Authorization: Bearer <token_do_aprendente>
Content-Type: application/json

{ "selectedOption": 1 }
\end{verbatim}

A resposta inclui \texttt{sessionResolved: Boolean}.
Quando \texttt{true}, o \texttt{DefaultStruggleComponent} emite
\texttt{StruggleEffect.NavigateToHome}, que o \texttt{DefaultRootComponent}
trata com \texttt{navigation.replaceAll(Config.Home)} — substituindo toda
a pilha para evitar que o aprendente regresse ao ecrã de dificuldade com
o botão de retrocesso.

\begin{figure}[htbp]
  \centering
  \begin{tikzpicture}[>=Stealth, font=\small,
    box/.style={draw, fill=#1, rounded corners=4pt,
                minimum width=4.2cm, minimum height=0.9cm, align=center},
    api/.style={draw, fill=gray!10, rounded corners=4pt,
                minimum width=4.2cm, minimum height=0.7cm, align=center}
  ]
    \node[box=red!10]    (s1) at (0, 0)    {StruggleScreen\\{\tiny lista de tarefas adaptativas}};
    \node[api]           (g1) at (5.8, 0)  {\texttt{GET /struggle/enrollment/\{id\}}};
    \node[box=orange!12] (s2) at (0,-2.0)  {Tarefa adaptativa 1\\{\tiny opções embaralhadas}};
    \node[api]           (g2) at (5.8,-2.0){\texttt{POST .../tasks/\{id\}/submit}};
    \node[box=orange!12] (s3) at (0,-3.5)  {Tarefa adaptativa N};
    \node[api]           (g3) at (5.8,-3.5){\texttt{POST .../tasks/\{id\}/submit}\\{\tiny sessionResolved=true}};
    \node[box=green!12]  (s4) at (0,-5.0)  {Home\\{\tiny replaceAll}};

    \draw[->] (s1) -- node[above,font=\tiny]{carregar} (g1);
    \draw[->] (g1.west) -- ++(-0.5,0) |- (s2.east);
    \draw[->] (s2) -- node[above,font=\tiny]{submeter} (g2);
    \draw[->] (g2.west) -- ++(-0.5,0) |- (s3.east);
    \draw[->] (s3) -- node[above,font=\tiny]{submeter} (g3);
    \draw[->] (g3.west) -- ++(-0.5,0) |- (s4.east);
  \end{tikzpicture}
  \caption{Fluxo de ecrã e pedidos HTTP durante uma sessão de remediação adaptativa}
  \label{fig:struggle_flow}
\end{figure}

% ─────────────────────────────────────────────────────────────────────────────
\section{Criação de Conteúdo: Engenharia de \textit{Prompts}}
\label{sec:engenharia_prompts}

Os \textit{prompts} do LLM do \textit{Play4Change} são elaborados segundo
o enquadramento de quatro componentes descrito pela Google no seu Guia de
\textit{Prompting} do Workspace~\cite{google2024promptguide}:
\textbf{Persona}, \textbf{Tarefa}, \textbf{Contexto} e \textbf{Formato}.
Cada componente está conscientemente presente em ambos os \textit{prompts}
utilizados pela plataforma.

\subsection*{\textit{Prompt} de geração de tarefas}

\paragraph{Persona (mensagem de sistema):}

\begin{verbatim}
You are an expert educational content designer
specialising in gamified learning experiences.
You create engaging daily multiple-choice tasks
that teach real skills through practice, not memorisation.
\end{verbatim}

Este enquadramento prepara o modelo para priorizar a aplicabilidade
prática em detrimento de factos de memorização, alinhando-se com a teoria
de aprendizagem construtivista~\cite{vygotsky1978mind} e com os princípios
de aprendizagem por mestria~\cite{bloom1984mastery}.

\paragraph{Formato (mensagem de sistema):}

\begin{verbatim}
OUTPUT FORMAT: Respond with a valid JSON array only.
No preamble. No explanation.
Schema:
[{
  "title": "string (max 60 chars)",
  "description": "string (max 300 chars)",
  "hint": "string (max 150 chars)",
  "pointsReward": number (10-100),
  "options": ["correta","errada1","errada2","errada3"],
  "correctAnswerIndex": 0
}]
\end{verbatim}

A imposição de JSON analisável automaticamente diretamente na instrução de
formato elimina a necessidade de heurísticas de pós-processamento na
maioria dos casos.
A convenção \texttt{correctAnswerIndex: 0} delega a mistura das opções ao
servidor, garantindo que o modelo nunca é responsável pela aleatorização.

\paragraph{Tarefa + Contexto (mensagem do utilizador):}

\begin{verbatim}
Generate 15 multiple-choice learning tasks for:

Subject Domain: Mobilidade Urbana
Module Objective: Compreender os princípios do transporte sustentável
Audience Level: Beginner — assumes no prior knowledge,
                uses simple language, concrete examples
Language: pt-PT

Important: 3 tasks already exist for this module.
Generate semantically distinct tasks that complement
the existing content.
\end{verbatim}

A variável de contexto \texttt{existingEmbeddings.size} alerta o modelo
para evitar sobreposição semântica, complementando a desduplicação por
cosseno executada após a geração.

\subsection*{\textit{Prompt} de análise de dificuldade}

O \textit{prompt} de dificuldade estende o mesmo enquadramento com contexto
de diagnóstico mais rico.
O papel de sistema muda de designer de conteúdo para treinador de
aprendizagem adaptativa:

\begin{verbatim}
You are an expert adaptive learning coach. A student
is struggling with a specific task in their learning
journey. Your job is to create a short sequence of
simpler subtasks that scaffold the student back to success.

PRINCIPLES:
- Diagnose the root cause from the error pattern
- Start with the simplest possible related concept
- Each subtask builds on the previous one
- The final subtask returns to the original task level
- Be encouraging — frame tasks as achievable steps
\end{verbatim}

A mensagem do utilizador fornece o contexto de diagnóstico completo,
incluindo o \texttt{ErrorPattern} classificado:

\begin{verbatim}
Diagnosed Struggle Type: The student misunderstands the
core concept. Start by re-explaining it differently.
\end{verbatim}

Os seis valores de \texttt{ErrorPattern} mapeiam cada um para uma
diretriz de diagnóstico distinta:

\begin{table}[ht]
  \caption{Diretivas de diagnóstico por padrão de erro no \texttt{StruggleAnalysisPrompt}}
  \label{tab:padroes_erro}
  \vskip0.5em
  \centering
  \begin{tabular}{lp{8cm}}
    \toprule
    \textbf{Padrão} & \textbf{Diretiva do \textit{prompt}} \\
    \midrule
    \texttt{CONCEPTUAL\_MISUNDERSTANDING} &
      Re-explicar o conceito central de forma diferente \\
    \texttt{PROCEDURAL\_ERROR} &
      Focar no procedimento / passos corretos \\
    \texttt{KNOWLEDGE\_GAP} &
      Preencher o conhecimento pré-requisito antes de regressar \\
    \texttt{UNCLEAR\_INSTRUCTIONS} &
      Fornecer orientação mais clara e explícita \\
    \texttt{TIME\_PRESSURE} &
      Usar tarefas de prática passo a passo ao ritmo do aprendente \\
    \texttt{UNKNOWN} &
      Fornecer andaimamento geral do básico ao intermédio \\
    \bottomrule
  \end{tabular}
\end{table}

\subsection*{Validação de esquema e nova tentativa}

Se a resposta do LLM tiver campos em falta ou JSON malformado,
\texttt{parseTasksFromJson()} devolve menos tarefas do que as solicitadas.
O adaptador acrescenta então \texttt{TaskGenerationPrompt.schemaReminder()}
à mensagem do utilizador e realiza uma única nova tentativa.
Se esta também falhar, o contador
\texttt{ai.generation.failures\{type="schema\_validation"\}} é
incrementado e é devolvido um \texttt{503}.
A estratégia de nova tentativa é conservadora (uma tentativa, sem
\textit{back-off} exponencial) porque a geração pelo Mistral é sensível à
latência e múltiplas tentativas bloqueariam a fila do conjunto de
\textit{threads}.

\subsection*{O volante de inércia da personalização}

A Figura~\ref{fig:volante_personalizacao} ilustra como o sistema acumula
inteligência pedagógica ao longo do tempo.

\begin{figure}[htbp]
  \centering
  \begin{tikzpicture}[>=Stealth, font=\small]
    \def\r{2.8}
    \draw[thick, ->] (0:\r) arc (0:350:\r);
    \node[draw, fill=blue!10, rounded corners=3pt, align=center,
          minimum width=2.2cm, minimum height=0.8cm]
          at (90:\r) {Aprendente tem\\dificuldade};
    \node[draw, fill=orange!10, rounded corners=3pt, align=center,
          minimum width=2.2cm, minimum height=0.8cm]
          at (0:\r) {Padrão de erro\\classificado};
    \node[draw, fill=purple!10, rounded corners=3pt, align=center,
          minimum width=2.2cm, minimum height=0.8cm]
          at (270:\r) {Ramo gerado\\ou reutilizado};
    \node[draw, fill=green!10, rounded corners=3pt, align=center,
          minimum width=2.2cm, minimum height=0.8cm]
          at (180:\r) {Ramo e \textit{embedding}\\persistidos};
    \node[align=center, font=\bfseries\small] at (0,0)
          {Volante de\\Personalização};
  \end{tikzpicture}
  \caption{O volante de inércia da personalização: cada dificuldade enriquece a reutilização futura}
  \label{fig:volante_personalizacao}
\end{figure}

Cada vez que um aprendente revela uma dificuldade, o ramo adaptativo
resultante é incorporado e armazenado.
Quando o aprendente seguinte encontra uma dificuldade semanticamente
similar — independentemente do tópico ou língua — o sistema recupera o
ramo anterior em vez de gerar um novo.
À medida que a base de utilizadores cresce, a taxa de reutilização aumenta,
reduzindo as chamadas ao LLM e a latência, e melhorando a consistência da
experiência adaptativa.
O arranque a frio (sem dados históricos) produz exclusivamente geração
inédita; o percurso quente emerge organicamente a partir do uso acumulado.

% ─────────────────────────────────────────────────────────────────────────────
\section{Observabilidade: Métricas, \textit{Dashboards} e Alertas}
\label{sec:observabilidade_deepdive}

\subsection*{O que significa observabilidade no \textit{Play4Change}}

A observabilidade é a capacidade de inferir o estado interno de um sistema
a partir das suas saídas externas~\cite{humble2018accelerate}.
No \textit{Play4Change}, cujo subsistema mais crítico é um agente de IA
com latência não determinística e modos de falha parcial, a observabilidade
não é um acrescento ao processo de implantação mas um requisito de desenho.
Três camadas de instrumentação são implementadas: \textbf{Micrometer}
(recolha de métricas JVM), \textbf{Prometheus} (armazenamento de séries
temporais e alertas) e \textbf{Grafana} (\textit{dashboards} e
visualização).

\subsection*{Métricas de domínio}

Para além dos valores predefinidos do Spring Boot Actuator (memória
\textit{heap} JVM, conjuntos de \textit{threads}, histogramas de latência
HTTP), são definidas sete métricas de domínio específicas:

\begin{table}[ht]
  \caption{Métricas Prometheus personalizadas no \textit{Play4Change}}
  \label{tab:metricas}
  \vskip0.5em
  \centering
  \begin{tabular}{llp{5.5cm}}
    \toprule
    \textbf{Métrica} & \textbf{Tipo} & \textbf{Descrição} \\
    \midrule
    \texttt{ai.generation.duration} & Temporizador &
      Duração da chamada ao LLM, rotulada por
      \texttt{generation\_phase} (GENERATION / ANALYSIS) \\
    \texttt{ai.generation.failures\_total} & Contador &
      Falhas do LLM por tipo
      (\texttt{schema\_validation}, \texttt{fixed\_path}, \texttt{adaptive}) \\
    \texttt{ai.adaptive.reuse} & Contador &
      Distribuição de estratégia de reutilização
      (\texttt{full}, \texttt{partial}, \texttt{fresh}) \\
    \texttt{ai.dedup.duplicates\_skipped} & Contador &
      Tarefas rejeitadas pela desduplicação por cosseno \\
    \texttt{struggle\_sessions\_total} & Contador &
      Novas sessões de dificuldade por \texttt{error\_pattern} \\
    \texttt{struggle\_sessions\_created\_total} & Contador &
      Sessões de dificuldade por tópico \\
    \texttt{topic\_enrollments\_total} & Contador &
      Inscrições por tópico \\
    \bottomrule
  \end{tabular}
\end{table}

\subsection*{Regras de alertas}

O Prometheus avalia oito regras de alerta definidas em
\texttt{alert.rules.yml}, abrangendo quatro níveis de severidade:

\begin{table}[ht]
  \caption{Regras de alerta Prometheus do \textit{Play4Change}}
  \label{tab:alertas}
  \vskip0.5em
  \centering
  \begin{tabular}{lll}
    \toprule
    \textbf{Alerta} & \textbf{Severidade} & \textbf{Condição} \\
    \midrule
    \texttt{ServerDown} & crítico &
      \texttt{up\{job="spring-boot"\} == 0} durante 1min \\
    \texttt{HighHttpErrorRate} & aviso &
      Taxa de erros 5xx $> 5\%$ durante 2min \\
    \texttt{HighAiGenerationFailureRate} & aviso &
      $> 0{,}1$ falhas/s durante 5min \\
    \texttt{AiCircuitBreakerOpen} & crítico &
      Estado do disjuntor aberto (imediato) \\
    \texttt{AiGenerationP95High} & aviso &
      Latência p95 $> 90$s durante 5min \\
    \texttt{JvmHeapUsageHigh} & aviso &
      \textit{heap} $> 85\%$ durante 5min \\
    \texttt{JvmHeapUsageCritical} & crítico &
      \textit{heap} $> 95\%$ durante 2min \\
    \texttt{GenerationThreadPoolQueueSaturated} & aviso &
      Fila $\geq 20$ tarefas durante 2min \\
    \bottomrule
  \end{tabular}
\end{table}

\subsection*{\textit{Dashboards} Grafana}

Três \textit{dashboards} Grafana são provisionados como JSON
(Infraestrutura como Código):

\begin{itemize}
  \item \textbf{Plataforma \textit{Play4Change}} — vista operacional de
    alto nível: taxa de pedidos HTTP, taxa de erros, inscrições ativas,
    \textit{badges} emitidos por hora.
  \item \textbf{Latência de IA} — durações das fases de geração,
    histogramas p50/p95/p99, distribuição da estratégia de reutilização,
    poupanças por desduplicação.
  \item \textbf{Fluxo do Aprendente} — funil de inscrições, taxa de
    sessões de dificuldade por tópico, evolução da reutilização de ramos
    adaptativos ao longo do tempo.
\end{itemize}

Os \textit{dashboards} são provisionados via API de provisionamento do
Grafana (ficheiros YAML em
\texttt{infra/grafana/provisioning/dashboards/}), sendo criados
automaticamente no arranque do contentor sem necessidade de importação
manual.
Esta abordagem garante que uma nova implantação é imediatamente observável
sem configuração adicional.

\begin{figure}[htbp]
  \centering
  \begin{tikzpicture}[>=Stealth, font=\small,
    comp/.style={draw, rounded corners=4pt, fill=gray!10,
                 minimum width=2.6cm, minimum height=0.9cm, align=center}
  ]
    \node[comp] (app)   at (0,0)    {Spring Boot\\Actuator};
    \node[comp] (prom)  at (4.2,0)  {Prometheus\\(scraping a 5s)};
    \node[comp] (graf)  at (8.4,0)  {Grafana\\Dashboards};
    \node[comp] (alert) at (4.2,-2.2){Alertmanager};

    \draw[->] (app)  -- node[above,font=\tiny]{/actuator/prometheus} (prom);
    \draw[->] (prom) -- node[above,font=\tiny]{consultas PromQL} (graf);
    \draw[->] (prom) -- node[left, font=\tiny]{alert.rules.yml} (alert);
    \draw[->, dashed] (alert) -- node[right,font=\tiny]{\textit{webhook}/e-mail} (4.2,-3.4);
  \end{tikzpicture}
  \caption{Pilha de observabilidade: Micrometer → Prometheus → Grafana}
  \label{fig:pilha_observabilidade}
\end{figure}

% ─────────────────────────────────────────────────────────────────────────────
\section{Arquitetura de Segurança}
\label{sec:seguranca_deepdive}

\subsection*{Segurança do transporte}

Todo o tráfego é encaminhado através do Nginx, que atua como
\textit{proxy} inverso com terminação TLS.
A configuração do Nginx impõe HTTPS e define os cabeçalhos
\texttt{X-Frame-Options: DENY}, \texttt{X-Content-Type-Options: nosniff}
e \texttt{Referrer-Policy: no-referrer}, mitigando ataques de
\textit{clickjacking}, deteção de MIME e fuga de \textit{referrer}.

\subsection*{Armazenamento de \textit{tokens}}

Os \textit{tokens} de acesso são armazenados em \texttt{sessionStorage},
que é isolado por separador do navegador e eliminado ao fechar o
separador, limitando a janela de exposição.
Os \textit{tokens} de atualização são armazenados num \textit{cookie} com
\texttt{SameSite=Strict; Secure; path=/}.
O atributo \texttt{SameSite=Strict} impede o envio do \textit{cookie}
em pedidos entre domínios, mitigando efetivamente CSRF sem a necessidade
de um \textit{token} CSRF separado~\cite{west2016samesite}.

\subsection*{Mapeamento OWASP Top 10}

\begin{table}[ht]
  \caption{Controlos do OWASP Top 10 (2021) implementados no \textit{Play4Change}}
  \label{tab:owasp}
  \vskip0.5em
  \centering
  \begin{tabular}{lp{8.5cm}}
    \toprule
    \textbf{Risco} & \textbf{Controlo implementado} \\
    \midrule
    A01 Controlo de Acesso Quebrado &
      \texttt{@AuthenticationPrincipal} vincula cada recurso ao principal
      autenticado; \textit{endpoints} de administração protegidos por
      verificações de papel do Spring Security \\
    A02 Falhas Criptográficas &
      SHA-256 para \textit{tokens} de \textit{magic link}; HTTPS/TLS para
      transporte; sem credenciais em texto simples armazenadas \\
    A03 Injeção &
      Todo o acesso à base de dados via JPA ou \texttt{JdbcTemplate}
      parametrizado; valores \textit{pgvector} convertidos via
      \texttt{?::vector} e não por concatenação de strings \\
    A05 Configuração de Segurança Incorreta &
      Cabeçalhos de segurança Nginx; CORS restrito a origens conhecidas;
      \textit{endpoints} do Spring Boot Actuator apenas na porta interna \\
    A07 Falhas de Identificação e Autenticação &
      \textit{Magic link} sem palavra-passe elimina reutilização de
      credenciais; consumo único do \textit{token}; limitação de taxa nos
      pedidos de ligação \\
    A09 Registo e Monitorização de Segurança &
      Registos estruturados SLF4J em cada ramo de lógica de negócio;
      alertas Prometheus sobre taxas de erro e falhas de autenticação \\
    \bottomrule
  \end{tabular}
\end{table}

\subsection*{Fuga de informação em respostas de erro}

O \textit{endpoint} de \textit{magic link} devolve a mesma resposta
\texttt{202 Accepted} independentemente de o endereço de correio eletrónico
ser conhecido ou não, impedindo enumeração de utilizadores.
Todas as respostas de erro de domínio seguem o esquema estruturado
\texttt{\{field, reason\}} sem expor ao cliente qualquer pilha de chamadas
ou estado interno.
As respostas \texttt{500 Internal Server Error} não contêm corpo em
ambiente de produção.

% ─────────────────────────────────────────────────────────────────────────────
\section{Tratamento de Erros: Uma Visão Sistémica}
\label{sec:tratamento_erros}

\subsection*{Tratamento funcional de erros com \textit{Arrow Either}}

O servidor utiliza o tipo de retorno \texttt{Either<AppError,~T>} da
biblioteca Arrow~\cite{arrow2024} em toda a camada de aplicação.
Cada subclasse de \texttt{AppError} carrega uma propriedade
\texttt{httpStatus} que o controlador utiliza para definir o código de
estado HTTP.
O compilador Kotlin obriga todos os locais de chamada a tratar ambos os
ramos, tornando os percursos de erro estruturalmente visíveis e impedindo
a propagação silenciosa.

\begin{table}[ht]
  \caption{Hierarquia de erros de domínio e códigos de estado HTTP}
  \label{tab:codigos_erro}
  \vskip0.5em
  \centering
  \begin{tabular}{lll}
    \toprule
    \textbf{Classe de erro} & \textbf{Estado HTTP} & \textbf{Causa típica} \\
    \midrule
    \texttt{BadRequest.InvalidField} & 400 &
      Falha de validação (campo em falta, valor inválido) \\
    \texttt{Unauthorized} & 401 &
      \textit{Token} de acesso ausente ou expirado \\
    \texttt{Forbidden} & 403 &
      Autenticado mas com papel insuficiente \\
    \texttt{NotFound.ResourceNotFound} & 404 &
      Entidade não existe \\
    \texttt{Conflict.ConcurrentModification} & 409 &
      Inscrição duplicada, regeneração concorrente \\
    \texttt{RateLimited} & 429 &
      Taxa de pedidos de \textit{magic link} excedida \\
    \texttt{InternalServerError} & 500 &
      Falha de carregamento para o MinIO, exceção inesperada \\
    \texttt{ServiceUnavailable.DependencyUnavailable} & 503 &
      API Mistral indisponível, falha de validação de esquema \\
    \bottomrule
  \end{tabular}
\end{table}

\subsection*{Modos de falha da IA}

A camada de IA introduz modos de falha ausentes nas APIs REST convencionais:

\begin{enumerate}
  \item \textbf{\textit{Timeout} de rede} — \texttt{withTimeoutOrNull(60\_000)}
    cancela a coroutine; a sessão de dificuldade é abandonada (estado
    \texttt{ABANDONED}).
  \item \textbf{Falha de validação de esquema} — \textit{array} JSON com
    campos obrigatórios em falta. Uma nova tentativa com
    \texttt{schemaReminder}; se persistir → \texttt{503}.
  \item \textbf{Discrepância de dimensão do \textit{embedding}} — alteração
    de versão da API do modelo devolve um vetor de tamanho diferente.
    \texttt{IllegalArgumentException} capturada → \texttt{503}.
  \item \textbf{Disjuntor aberto} — Resilience4j acciona o disjuntor após
    falhas repetidas; as chamadas subsequentes são curto-circuitadas
    imediatamente → \texttt{503}.
    O alerta Prometheus \texttt{AiCircuitBreakerOpen} é disparado.
  \item \textbf{Geração parcial} — algumas tarefas são analisadas
    corretamente, outras falham. As tarefas válidas são persistidas; se o
    total ficar abaixo do solicitado, o tópico transita na mesma para
    \texttt{ACTIVE} com um número reduzido de tarefas (degradação
    aceitável).
\end{enumerate}

\subsection*{Tratamento de erros no portal web}

No \textit{frontend} React, \texttt{useCreateTopicFromUrl} e
\texttt{useCreateTopicFromPdf} são mutações do TanStack Query.
O sinalizador \texttt{isError} activa uma mensagem de erro \textit{inline}
que extrai o \texttt{\{field, reason\}} estruturado da resposta de erro
Axios quando disponível, ou recorre a uma cadeia de caracteres genérica
localizada.
O componente \texttt{ErrorBoundary} captura erros de renderização não
tratados e apresenta uma interface de recuperação, impedindo uma página em
branco em caso de falha inesperada.

\subsection*{Tratamento de erros na aplicação móvel}

No cliente móvel, os erros seguem uma pipeline de transformação em três
etapas antes de chegarem à camada de apresentação
(Figura~\ref{fig:mobile_erros}):

\begin{enumerate}
  \item \textbf{Camada de rede → \texttt{NetworkError}} —
    O \texttt{NetworkErrorMapper} transforma o código de estado HTTP ou a
    exceção Ktor no tipo selado \texttt{NetworkError}
    (e.g., \texttt{401} → \texttt{Unauthorized},
    \texttt{ConnectTimeoutException} → \texttt{Timeout},
    \texttt{409} com cabeçalho de dificuldade → \texttt{StruggleOpen}).

  \item \textbf{\texttt{NetworkError} → \texttt{AppError}} —
    O método de extensão \texttt{NetworkError.toAppError()} converte para
    o tipo selado de domínio \texttt{AppError}, independente de HTTP.
    Este mapeamento mantém a camada de apresentação agnóstica em relação ao
    protocolo de transporte.

  \item \textbf{\texttt{AppError} → mensagem localizada} —
    O \texttt{AppError.messageKey} é resolvido para a cadeia de
    caracteres localizada pelo sistema de recursos do Compose Multiplatform.
    O \texttt{BaseView} renderiza um \texttt{ErrorDialog} automaticamente
    quando \texttt{state.error != null}, sem necessidade de lógica
    adicional nos ecrãs individuais.
\end{enumerate}

\begin{figure}[htbp]
  \centering
  \begin{tikzpicture}[>=Stealth, font=\small,
    box/.style={draw, fill=#1, rounded corners=4pt,
                minimum width=3.5cm, minimum height=0.9cm, align=center}
  ]
    \node[box=red!10]    (net) at (0,0)    {Ktor / Exceção\\HTTP};
    \node[box=orange!12] (ne)  at (4.2,0)  {\texttt{NetworkError}\\(selado, 8 casos)};
    \node[box=yellow!18] (ae)  at (8.4,0)  {\texttt{AppError}\\(selado, agnóstico)};
    \node[box=green!12]  (ui)  at (8.4,-2.2){\texttt{ErrorDialog}\\+ \texttt{messageKey}};

    \draw[->] (net) -- node[above,font=\tiny]{\texttt{toNetworkError()}} (ne);
    \draw[->] (ne)  -- node[above,font=\tiny]{\texttt{toAppError()}} (ae);
    \draw[->] (ae)  -- node[right,font=\tiny]{\texttt{BaseView}} (ui);
  \end{tikzpicture}
  \caption{Pipeline de transformação de erros no cliente móvel KMP}
  \label{fig:mobile_erros}
\end{figure}

% ─────────────────────────────────────────────────────────────────────────────
\section{Superfície Completa da API REST}
\label{sec:api_completa}

\begin{table}[ht]
  \caption{Superfície completa da API REST do \textit{Play4Change}}
  \label{tab:api_completa}
  \vskip0.5em
  \centering
  \footnotesize
  \begin{tabular}{lll}
    \toprule
    \textbf{Método e Caminho} & \textbf{Ator} & \textbf{Finalidade} \\
    \midrule
    \texttt{POST /auth/magic-link}               & Qualquer & Solicitar \textit{magic link} \\
    \texttt{GET  /auth/verify}                   & Qualquer & Redirecionamento para \textit{deep link} móvel \\
    \texttt{POST /auth/magic-link/verify}        & Qualquer & Verificar \textit{token} → par JWT \\
    \texttt{POST /auth/refresh}                  & Qualquer & Renovar \textit{token} de acesso \\
    \texttt{POST /auth/logout}                   & Qualquer & Revogar \textit{token} de atualização \\
    \midrule
    \texttt{GET  /admin/me}                      & Admin. & Perfil do administrador atual \\
    \texttt{POST /admin/topics}                  & Admin. & Criar tópico a partir de URL \\
    \texttt{POST /admin/topics/pdf}              & Admin. & Criar tópico a partir de PDF \\
    \texttt{GET  /admin/topics}                  & Admin. & Listar tópicos (paginado) \\
    \texttt{GET  /admin/topics/\{id\}}           & Admin. & Detalhes do tópico e registo de fases \\
    \texttt{GET  /admin/topics/\{id\}/progress}  & Admin. & Fluxo SSE de geração \\
    \texttt{POST /admin/topics/\{id\}/regenerate}& Admin. & Reiniciar \textit{pipeline} de geração \\
    \texttt{GET  /admin/topics/\{id\}/tasks}     & Admin. & Listar modelos de tarefas \\
    \texttt{PUT  /admin/tasks/\{id\}}            & Admin. & Editar modelo de tarefa \\
    \texttt{GET  /admin/topics/\{id\}/struggle-tasks} & Admin. & Listar tarefas adaptativas \\
    \texttt{PUT  /admin/adaptive-tasks/\{id\}}   & Admin. & Editar tarefa adaptativa \\
    \texttt{GET  /admin/topics/\{id\}/struggle-path-stats} & Admin. & Estatísticas de dificuldade \\
    \texttt{GET  /admin/topics/\{id\}/prerequisites}& Admin. & Obter pré-requisitos \\
    \texttt{POST /admin/topics/\{id\}/prerequisites}& Admin. & Definir pré-requisitos \\
    \texttt{GET  /admin/learning-graph}          & Admin. & DAG completo de pré-requisitos \\
    \texttt{GET  /admin/topics/\{id\}/badges}    & Admin. & Estatísticas de \textit{badges} por tópico \\
    \texttt{GET  /admin/users}                   & Admin. & Listar utilizadores (paginado) \\
    \texttt{GET  /admin/users/\{id\}}            & Admin. & Detalhes e estatísticas do utilizador \\
    \texttt{POST /admin/users/\{id\}/promote}    & Admin. & Promover para papel de administrador \\
    \texttt{GET  /admin/users/\{id\}/enrollments}& Admin. & Lista de inscrições do utilizador \\
    \texttt{GET  /admin/users/\{id\}/badges}     & Admin. & Lista de \textit{badges} do utilizador \\
    \texttt{GET  /admin/users/\{id\}/enrollments/\{eid\}/roadmap} & Admin. & Mapa de aprendizagem \\
    \midrule
    \texttt{POST /user/enroll}                   & Aprendente & Inscrever em tópico \\
    \texttt{POST /user/tasks/\{id\}/submit}      & Aprendente & Submeter resposta a tarefa \\
    \texttt{GET  /user/roadmap}                  & Aprendente & Mapa de aprendizagem do aprendente \\
    \texttt{GET  /user/topics}                   & Aprendente & Tópicos disponíveis \\
    \texttt{GET  /user/struggle/\{sessionId\}}   & Aprendente & Tarefas do ramo adaptativo \\
    \texttt{POST /user/struggle/\{taskId\}/submit}& Aprendente & Submeter tarefa adaptativa \\
    \texttt{GET  /user/badges}                   & Aprendente & \textit{Badges} conquistados \\
    \texttt{GET  /user/profile}                  & Aprendente & Perfil e preferências \\
    \texttt{PUT  /user/profile/name}             & Aprendente & Atualizar nome de apresentação \\
    \texttt{GET  /public/stats}                  & Público & Estatísticas da página de entrada \\
    \bottomrule
  \end{tabular}
\end{table}

\paragraph{Referências arquiteturais.}
A arquitetura do servidor segue a Arquitetura Hexagonal (Portas e
Adaptadores)~\cite{cockburn2005hexagonal} e o
\textit{Domain-Driven Design}~\cite{evans2003ddd}.
O tratamento funcional de erros adopta o padrão \textit{Either} da
biblioteca Arrow~\cite{arrow2024}.
O \textit{streaming} SSE reactivo assenta na especificação W3C de
\textit{Server-Sent Events}~\cite{w3c2021sse}.
A pesquisa por similaridade vectorial utiliza a extensão PostgreSQL
\textit{pgvector}~\cite{pgvector2024} com indexação aproximada de
vizinho mais próximo pelo algoritmo IVFFlat~\cite{johnson2019billion}.
A engenharia de \textit{prompts} segue o enquadramento de quatro
componentes do Guia de \textit{Prompting} da Google~\cite{google2024promptguide}
(Persona, Tarefa, Contexto, Formato).
A estratégia de ramos adaptativos assenta na teoria do andaimamento
pedagógico~\cite{vygotsky1978mind} e nos princípios de aprendizagem
por mestria~\cite{bloom1984mastery}.
